\documentclass[11pt,a4paper]{article}
\usepackage{DDTA_METHODOLOGY_GUIDE_STYLE_R1}

\hypersetup{
  pdftitle={DDTA Documentation Authoring Guide - Revision 7 Rebuild R8},
  pdfauthor={Documentation-Driven Threat Analysis research project}
}

% -----------------------------------------------------------------------------
% PAGE-INTEGRITY HASHES
% Generated by update_page_md5.py over canonical LaTeX payloads between
% PAGE-CONTENT markers. Visible self-hash lines are excluded. The integrity
% index resolves hashes of all listed pages and canonicalizes its own hash as
% <SELF> to avoid circular self-reference.
% -----------------------------------------------------------------------------
\newcommand{\PageMDfive}[1]{\csname PageMDfive#1\endcsname}
\providecommand{\tightlist}{\setlength{\itemsep}{0pt}\setlength{\parskip}{0pt}}
\providecommand{\passthrough}[1]{#1}
%<PAGE-HASH-DEFS>
\expandafter\def\csname PageMDfive1\endcsname{7289db2d0692a47b1d1013b12077997f}
\expandafter\def\csname PageMDfive2\endcsname{14ee0b08318ee0c7e0fbe0df6796b7e5}
\expandafter\def\csname PageMDfive3\endcsname{22e57dfc4b4f586235f8aed3c522a7e2}
\expandafter\def\csname PageMDfive4\endcsname{3bdac4b5fec97501c23d7968bf1619d5}
\expandafter\def\csname PageMDfive5\endcsname{b69ae801bbd9fab0dd9722c5416054c4}
\expandafter\def\csname PageMDfive6\endcsname{dd6326ef067e8650d8c2fb33662caf1b}
\expandafter\def\csname PageMDfive7\endcsname{8828e008581efc3d40954e198bdb8c0c}
\expandafter\def\csname PageMDfive8\endcsname{39dd4da3676b9d5a04cebdd9a196c627}
\expandafter\def\csname PageMDfive9\endcsname{ddcc227fc8f2ef02b4235ea1a2979c66}
\expandafter\def\csname PageMDfive10\endcsname{fee9637e13dfb1e7c117390773a8e3ae}
\expandafter\def\csname PageMDfive11\endcsname{97e6e0faea39d999309beffb6aca06bd}
\expandafter\def\csname PageMDfive12\endcsname{70297884e3aacfc8f40b330ff7183ca4}
\expandafter\def\csname PageMDfive13\endcsname{9a97d1490def76a4e8d0cd67a8307a34}
\expandafter\def\csname PageMDfive14\endcsname{b008d2293cbe637ae1268cc43c3b2351}
\expandafter\def\csname PageMDfive15\endcsname{af9687992a8bac51cb56ea985aeb2ebc}
\expandafter\def\csname PageMDfive16\endcsname{fab634987bb3e565df9c1d1f7ad066d0}
\expandafter\def\csname PageMDfive17\endcsname{bcb1bb076c245425bda75ebd8a2cb5c6}
\expandafter\def\csname PageMDfive18\endcsname{a70c8156b1d0d0ddd6c5640bdb2a8564}
\expandafter\def\csname PageMDfive19\endcsname{62e667a1c7db0a35813bdb36a8115714}
\expandafter\def\csname PageMDfive20\endcsname{d240cc1ea1cb104187dcd2ce3d89b046}
\expandafter\def\csname PageMDfive21\endcsname{da7ba071ce540507ac258349820dd96d}
\expandafter\def\csname PageMDfive22\endcsname{a70410d72d2daf48666b8f2addac2407}
\expandafter\def\csname PageMDfive23\endcsname{2aed4a907b4ecca2027ecc165ad9dacc}
\expandafter\def\csname PageMDfive24\endcsname{df4e6b117b67e0101047ac7b62483340}
\expandafter\def\csname PageMDfive25\endcsname{8a63264d93dc6926960764a7b017eb86}
\expandafter\def\csname PageMDfive26\endcsname{e03bfb85cd279ff18b7e7b73a2028473}
\expandafter\def\csname PageMDfive27\endcsname{1f804382f39d12572469d7fb1269fd71}
\expandafter\def\csname PageMDfive28\endcsname{651725f978e9d26d4bf4afc05c31f933}
\expandafter\def\csname PageMDfive29\endcsname{052bb1e5051941bc7f21945436cffa4e}
\expandafter\def\csname PageMDfive30\endcsname{f5310549d08de40525cf85c2df1b4e16}
\expandafter\def\csname PageMDfive31\endcsname{8921ecd8cfa9bf31d7449e9c7da8ebbd}
\expandafter\def\csname PageMDfive32\endcsname{2565613776f5cf4b81ef70950950637d}
\expandafter\def\csname PageMDfive33\endcsname{2541b4df42daadc5fb2402472d4e8c10}
\expandafter\def\csname PageMDfive34\endcsname{255ec7ee9d4053db3e78fe65baaf4216}
\expandafter\def\csname PageMDfive35\endcsname{ddda2e275d829cad56c4dadb9a70d4ae}
\expandafter\def\csname PageMDfive36\endcsname{f15a3bb4200cf549e5828495b75c4384}
\expandafter\def\csname PageMDfive37\endcsname{e4a9d2dd8df5e956a76592038b883979}
\expandafter\def\csname PageMDfive38\endcsname{75fb1fdc83af88cce5cf926df26987e1}
\expandafter\def\csname PageMDfive39\endcsname{76fa920eae92f2bb829a2743a84b30a8}
\expandafter\def\csname PageMDfive40\endcsname{e7136268d32bf6086596d5f030a947ad}
\expandafter\def\csname PageMDfive41\endcsname{9bc0e63e85d2d026b925a81ded38de52}
\expandafter\def\csname PageMDfive42\endcsname{cc250cd2a2deb3069dd6892e943bd284}
\expandafter\def\csname PageMDfive43\endcsname{9c03ac3e778c27333ef2d51de02f6eb2}
\expandafter\def\csname PageMDfive44\endcsname{fbc67bc514bb10a6d2c69b4728ffd1ba}
\expandafter\def\csname PageMDfive45\endcsname{e7298e5d8dcc77d5ea99741209730c0c}
\expandafter\def\csname PageMDfive46\endcsname{477aa90c4377838130c76cffe2e7a58d}
%</PAGE-HASH-DEFS>

\newcommand{\PageHashFooter}[1]{%
  \fancyfoot[C]{}%
  \fancyfoot[R]{\sffamily\scriptsize\color{DDTAGray}Pag. \thepage\quad MD5 contenuto: \texttt{#1}}%
}

\newcommand{\IntegrityRow}[3]{#1 & #2 & {\footnotesize\ttfamily #3} \\
}

\begin{document}
\selectlanguage{italian}
\fancyhead[L]{\sffamily\small\color{DDTAGray}DDTA Documentation Authoring Guide}
\fancyhead[R]{\sffamily\small\color{DDTABlue}Revision 7 - Rebuild R8}
\fancyfoot[L]{\sffamily\scriptsize\color{DDTAGray}Documentation authoring methodology}

%<PAGE-DESC:1>Frontespizio della nuova guida di authoring
%<PAGE-CONTENT:1>
\begin{titlepage}
\thispagestyle{empty}
\vspace*{12mm}
{\sffamily\bfseries\color{DDTANavy}\fontsize{15}{18}\selectfont Documentation-Driven Threat Analysis}\par
\vspace{5mm}
{\sffamily\bfseries\color{DDTANavy}\fontsize{28}{32}\selectfont Documentation Authoring Guide}\par
\vspace{2mm}
{\sffamily\color{DDTABlue}\fontsize{18}{22}\selectfont Revision 7 - Rebuild R8}\par

\vspace{15mm}
\begin{tcolorbox}[enhanced,arc=3pt,boxrule=0pt,colback=DDTANavy,left=13pt,right=13pt,top=12pt,bottom=12pt]
{\color{white}\sffamily\bfseries\large Nuova costruzione della guida DDTA per la scrittura della documentazione}\par
\vspace{2mm}
{\color{white!88}\small Il documento viene ricostruito progressivamente da una baseline vuota. Ogni nuova sezione sara' aggiunta solo dopo revisione del significato metodologico e verifica dell'integrita' delle pagine gia' approvate.}
\end{tcolorbox}

\vfill
{\sffamily\scriptsize\color{DDTAGray}MD5 contenuto pagina 1: \texttt{\PageMDfive{1}}}\par
\end{titlepage}
%</PAGE-CONTENT:1>

\clearpage
\setcounter{page}{2}
%<PAGE-DESC:2>Project Problem Framing - definizione e regole di authoring
%<PAGE-CONTENT:2>
\PageHashFooter{\PageMDfive{2}}
\section{Step 0 - Project Problem Framing}

\begin{ddtaopen}[title={RIQUADRO TEMPORANEO - WORKING SET E RIFERIMENTI - DA CANCELLARE PRIMA DELLA VERSIONE FINALE}]
{\scriptsize
\textbf{Artefatti attivi:} \nolinkurl{DDTA_DOCUMENTATION_AUTHORING_GUIDE_R7_REBUILD_R8.tex}; \nolinkurl{DDTA_BASE_ANALYSIS_GUIDE_REBUILD_R7.tex}; \nolinkurl{DDTA_DERMATRIAGE_PARALLEL_CASE_STUDY_R25_FR03_SPLIT_REVIEW_R1.tex}.\par
\textbf{Controllo di lavoro:} \nolinkurl{DDTA_R25_BASE_ANALYSIS_GUIDE_REBUILD_WORK_PLAN_R6.md}; supporto di analisi: \nolinkurl{DDTA_R25_DERMATRIAGE_CASE_STUDY_WORKING_ANALYSIS_R1.md}.\par
\textbf{Riferimenti:} \nolinkurl{DDTA_DOCUMENTATION_AUTHORING_GUIDE_R6_CANDIDATE_R2.tex}; \nolinkurl{DDTA_BASE_ANALYSIS_OPERATIONAL_GUIDE_R3.tex} + BA0--BA5; \nolinkurl{DDTA_R25_BASE_ANALYSIS_REBUILD_SOURCE_REGISTER_R1.md}; Core/Complete BA candidate R1 come research evidence.\par
\textbf{Disciplina corrente:} il piano R6 apre prima un audit completo di preservazione semantica delle sei fonti originali DermaTriage rispetto alla documentazione DDTA corrente. Ogni perdita o indebolimento viene corretto nel layer documentale prima di rieseguire la BA; i gap della fonte restano espliciti. Solo dopo questo gate si esegue l'audit di coerenza/connettivita' dell'intera BA accettata e, successivamente, la review dei costrutti ancora aperti.\par
}
\end{ddtaopen}

\subsection{Che cos'e' il framing}
Il \textbf{Project Problem Framing} mette a fuoco il problema --- o la classe di problemi --- che il progetto deve affrontare e il contributo principale che intende portare entro un boundary governato. Non e' un tipo documentale L1: resta una \textbf{precondizione dell'authoring DDTA} e non sostituisce MacroRequirement, Decision o FunctionalRequirement. Il boundary puo' contenere responsabilita' che verranno rese esplicite solo nei MacroRequirement successivi. Stare dentro il boundary non significa che ogni responsabilita', attore, integrazione o fase del ciclo di vita debba comparire nella frase finale di framing. Il framing deve rendere riconoscibile il problema; non deve anticipare il catalogo delle responsabilita' del progetto.

\begin{ddtaprinciple}[title={Domanda fondamentale}]
Quale problema o classe di problemi deve affrontare il progetto, entro quale boundary significativo, se rimuoviamo temporaneamente le scelte di soluzione correnti?
\end{ddtaprinciple}

\subsection{Regole di authoring}
Quello che va scritto e' il problema, non la pipeline o il modello scelto per risolverlo. Si puo' indicare il risultato atteso a livello di problema, ma senza trasformarlo in una sequenza di requisiti. Le regole operative, per esempio una soglia o un fallback, restano fuori: appartengono a livelli successivi. Ogni significato aggiunto deve essere supportato dalle fonti, e aggiungere dettaglio per riempire i campi non serve.

\begin{ddtacore}[title={Regola di neutralizzazione della soluzione}]
Per riconoscere il problema si mettono \emph{temporaneamente} da parte le scelte tecniche correnti. Questo non autorizza a perdere informazione: le scelte realmente governate verranno recuperate al semantic owner corretto nei livelli successivi.
\end{ddtacore}
%</PAGE-CONTENT:2>

\clearpage
%<PAGE-DESC:3>Project Problem Framing - procedura umana, gate ed esempio DermaTriage
%<PAGE-CONTENT:3>
\PageHashFooter{\PageMDfive{3}}
\subsection{Procedura umana di costruzione}

\begin{enumerate}
  \item \textbf{Leggere le fonti autorizzate prima di scrivere.} Identificare frasi che descrivono scopo, problema, outcome di dominio e boundary, senza assumere che ogni dettaglio sorgente appartenga al framing.
  \item \textbf{Separare problema e soluzione.} Mettere temporaneamente da parte i dettagli di realizzazione, per esempio una pipeline o una soglia, come possibili contenuti di livello successivo.
  \item \textbf{Formulare il problema in linguaggio di dominio.} La frase deve restare vera anche se le scelte tecniche correnti vengono sostituite, salvo che una specifica tecnologia costituisca essa stessa un vincolo governato del problema.
  \item \textbf{Controllare il boundary senza trasformarlo nel framing.} Non ampliare il progetto con obiettivi, attori o responsabilita' che le fonti non autorizzano. Allo stesso tempo, non copiare automaticamente nel framing finale tutto cio' che appartiene al boundary: una responsabilita' puo' essere interna al progetto e restare correttamente rinviata al livello MacroRequirement.
  \item \textbf{Fermarsi quando il problema e' riconoscibile.} MacroRequirement, Decision e FunctionalRequirement arrivano dopo.
\end{enumerate}

\begin{ddtacheck}[title={Gate di accettazione del framing}]
Prima di accettare il testo verificare:
\begin{itemize}
  \item se tutte le scelte tecniche correnti cambiassero, il problema descritto continuerebbe a esistere?
  \item il testo descrive un problema/contributo di progetto oppure sta gia' prescrivendo una soluzione?
  \item sono stati introdotti significati non supportati dalle fonti autorizzate?
  \item ci sono dettagli che possono essere rimandati senza perdere l'identita' del problema?
\end{itemize}
Se una risposta segnala contaminazione o informazione non autorizzata, occorre riformulare prima di procedere ai MacroRequirement.
\end{ddtacheck}

\begin{ddtaanti}[title={Anti-pattern principali}]
\begin{itemize}
  \item trasformare l'architettura corrente nella definizione del problema;
  \item elencare dettagli tecnici solo per rendere il framing piu' concreto;
  \item anticipare fallback o regole condizionali per spiegare casi speciali;
  \item copiare lo stesso significato che verra' poi riscritto integralmente nel primo MacroRequirement;
  \item introdurre obiettivi plausibili di dominio che le fonti non governano.
\end{itemize}
\end{ddtaanti}

\begin{ddtaexample}[title={DermaTriage - framing ottenuto durante l'applicazione della guida}]
\emph{Il problema che DermaTriage deve affrontare e' il supporto al triage precoce di casi dermatologici relativi a lesioni potenzialmente oncologiche. Il progetto parte dalle informazioni gia' disponibili sul caso per determinarne l'urgenza e una priorita' operativa di presa in carico, con l'obiettivo di favorire un instradamento specialistico tempestivo.}
\end{ddtaexample}

Nell'esempio non compaiono dettagli di realizzazione come la pipeline a quattro stadi o il fallback senza immagine. Queste informazioni non vengono eliminate dal progetto: vengono rinviate al livello documentale che ne possiede il significato.
%</PAGE-CONTENT:3>

\clearpage
%<PAGE-DESC:4>Project Problem Framing - LLM Execution Profile e prompt canonico
%<PAGE-CONTENT:4>
\PageHashFooter{\PageMDfive{4}}
\subsection{Profilo di esecuzione LLM}

Il profilo LLM adatta la guida umana all'esecuzione automatica. Non e' una metodologia alternativa e non estende l'authority DDTA. Il suo scopo e' contenere alcune tendenze ricorrenti dei modelli linguistici: completare oltre il necessario, anticipare la soluzione, riempire i campi per inerzia e spacciare una parafrasi per decomposizione.

\subsubsection{Prompt canonico - Project Problem Framing}
\begin{lstlisting}[basicstyle=\ttfamily\footnotesize]
ROLE

You are helping with DDTA documentation authoring. The DDTA Documentation
Authoring Guide is the methodological authority. Work only from the authorized
project sources provided for this step. Do not import project truth from domain
knowledge, earlier reconstructions, Base Analysis, threat models, tooling
requirements, or implementation assumptions.

TASK

Derive the Project Problem Framing.

Framing is a precondition for authoring, not an L1 document type. Identify the
problem or class of problems the project must address, together with the
contribution it intends to make, while setting aside current solution choices.

SOURCE-FIRST RULE

1. Pull out statements supported by the sources that describe purpose, problem,
   domain outcome, actors, and meaningful boundary.
2. Separate problem meaning from current realization details.
3. Keep realization details, but mark them DEFERRED so they can be placed at
   the right semantic level later.

WHAT NOT TO PUT IN THE FINAL FRAMING

- names of components, models, algorithms, or protocols just because they exist;
- pipeline stages or architecture choices;
- thresholds, mappings, fallback rules, or interface details;
- responsibilities or goals the authorized sources do not support;
- details added only to make the framing look complete.

COUNTERFACTUAL CHECK

Ask: "If all current technical and architectural choices were replaced, would
 the stated project problem still exist?" If not, rework the framing unless the
relevant choice is itself a governed problem boundary.

STOP RULE

Stop once the project problem and contribution are recognizable. Do not
pre-write MacroRequirements, Decisions, or FunctionalRequirements.

OUTPUT

Return exactly the analysis schema and final-output schema defined in this
section. Do not omit required fields. For REQUIRED-SLOT fields with no supported
value, use "-".
\end{lstlisting}
%</PAGE-CONTENT:4>

\clearpage
%<PAGE-DESC:5>Project Problem Framing - regola LLM e classificazione dei campi
%<PAGE-CONTENT:5>
\PageHashFooter{\PageMDfive{5}}
\begin{ddtaanti}[title={Regola per l'LLM}]
Un'informazione presente nella source non appartiene automaticamente al framing. Il modello deve prima classificarne il semantic level: conoscere un dettaglio non e' una ragione sufficiente per anticiparlo.\par
Anche un significato correttamente incluso nel project boundary non appartiene automaticamente alla frase finale di framing. Se puo' essere rinviato ai MacroRequirement senza perdere l'identita' del problema, deve restare fuori dal framing finale.
\end{ddtaanti}

\subsection{Schema di output per l'LLM}

Per il framing si tengono separati due livelli: la \textbf{worksheet di analisi} e l'\textbf{output documentale finale}. La worksheet serve a rendere verificabile il ragionamento del modello; non va copiata nella documentazione governata.

\subsubsection{Classificazione dei campi}
\begin{tabularx}{\textwidth}{L{48mm}L{34mm}Y}
\toprule
\textbf{Campo} & \textbf{Classe} & \textbf{Regola} \\
\midrule
Authorized sources & REQUIRED-CONTENT & Elenco delle sole fonti autorizzate per il passo. \\
Source-supported problem meaning & REQUIRED-CONTENT & Fatti di problema direttamente sostenuti dalle fonti. \\
Deferred solution details & REQUIRED-SLOT & Dettagli riconosciuti ma rinviati; ``-'' se assenti. \\
Unsupported / ambiguous claims & REQUIRED-SLOT & Significati non autorizzati o ancora ambigui; ``-'' se assenti. \\
Meaningful boundary & REQUIRED-SLOT & Boundary governato; ``-'' se non stabilito separatamente. \\
Counterfactual result & REQUIRED-CONTENT & PASS o REWORK con motivazione breve. \\
Final framing & REQUIRED-CONTENT & Un blocco conciso di prosa governata. \\
\bottomrule
\end{tabularx}
%</PAGE-CONTENT:5>

\clearpage
%<PAGE-DESC:6>Project Problem Framing - formato esatto dell'output LLM
%<PAGE-CONTENT:6>
\PageHashFooter{\PageMDfive{6}}
\subsubsection{Formato esatto richiesto}
\begin{lstlisting}[basicstyle=\ttfamily\footnotesize]
[DDTA FRAMING ANALYSIS]
AUTHORIZED SOURCES:
- ...

SOURCE-SUPPORTED PROBLEM MEANING:
- ...

MEANINGFUL BOUNDARY:
-

DEFERRED SOLUTION DETAILS:
- ...

UNSUPPORTED / AMBIGUOUS CLAIMS:
- ...

COUNTERFACTUAL CHECK:
PASS | REWORK
Rationale: ...

FRAMING VERDICT:
PROCEED | REWORK | STOP
Rationale: ...

[FINAL PROJECT PROBLEM FRAMING]
<one concise governed prose block>
\end{lstlisting}
%</PAGE-CONTENT:6>

\clearpage
%<PAGE-DESC:7>Project Problem Framing - gate LLM e nota sui controlli quantitativi
%<PAGE-CONTENT:7>
\PageHashFooter{\PageMDfive{7}}
\subsection{Gate sull'output LLM}

\begin{ddtacheck}[title={Controlli prima di accettare l'output LLM}]
\begin{itemize}
  \item il blocco finale contiene solo significato gia' verificato nella worksheet;
  \item i dettagli marcati DEFERRED non ricompaiono nel framing finale senza una ragione esplicita;
  \item nessun campo vuoto viene riempito con una parafrasi di un altro campo;
  \item \texttt{PROCEED} e' ammesso solo se il counterfactual check e la source-authority review sono superati;
  \item l'output LLM resta una proposta da sottoporre alla stessa review semantica prevista per l'autore umano.
\end{itemize}
\end{ddtacheck}

\begin{ddtacore}[title={Nota sui controlli quantitativi}]
Le metriche consolidate di similarita' e retrieval del corpus (Jaccard su shingles, cosine TF--IDF, cosine su embeddings e BM25 per retrieval) verranno definite in una sezione metodologica dedicata. Non sostituiscono i gate semantici e non sono necessarie per accettare il primo framing di un corpus privo di altri elementi documentali.
\end{ddtacore}

\begin{ddtaprinciple}[title={Separazione tra diagnostica e decisione metodologica}]
Le metriche quantitative potranno segnalare sovrapposizioni o possibili rewording, ma non potranno decidere da sole se un elemento DDTA e' corretto. La decisione resta affidata ai gate semantici previsti dalla guida.
\end{ddtaprinciple}
%</PAGE-CONTENT:7>

\clearpage
%<PAGE-DESC:8>Uso della metodologia - authoring nativo e reconstruction
%<PAGE-CONTENT:8>
\PageHashFooter{\PageMDfive{8}}
\section{Uso della metodologia: authoring nativo e reconstruction}

DDTA nasce come metodologia per l'\textbf{authoring nativo della documentazione di progetto}. Il suo impiego naturale comincia insieme al progetto e ne accompagna la progressiva definizione: problema, responsabilita', decisioni, requisiti.

Ricostruire documentazione preesistente e' invece un uso secondario, di validazione: serve a recuperare il significato governabile dalle fonti disponibili e a far emergere ambiguita', lacune e sovrapposizioni che una documentazione DDTA nativa avrebbe dovuto prevenire o quantomeno rendere esplicite.

\begin{ddtacore}[title={Tenere separati i due percorsi}]
Le regole per leggere documentazione legacy non devono appesantire il percorso di authoring nativo.
\end{ddtacore}

\begin{tabularx}{\textwidth}{L{36mm}Y}
\toprule
\textbf{Percorso} & \textbf{Punto di partenza e uso} \\
\midrule
\textbf{NATIVE} & Si parte dal problema governato e dalle responsabilita' che il progetto deve assumere. Le scelte di soluzione vengono introdotte ai livelli successivi quando diventano commitment governati. \\
\textbf{RECONSTRUCTION} & Si parte da fonti legacy autorizzate e si recupera soltanto il significato che esse supportano. Ambiguita', conflitti e significati mancanti restano visibili invece di essere completati per inferenza. \\
\bottomrule
\end{tabularx}

\begin{ddtaprinciple}[title={Stessa semantica, diverso punto di ingresso}]
I costrutti DDTA non cambiano tra i due percorsi. Cambia il modo in cui si ottiene il significato da documentare: nel native authoring viene governato mentre il progetto prende forma; nella reconstruction viene recuperato da una documentazione che non era stata scritta secondo DDTA.
\end{ddtaprinciple}
%</PAGE-CONTENT:8>

\clearpage
%<PAGE-DESC:9>MacroRequirement - definizione e rapporto con framing e Decision
%<PAGE-CONTENT:9>
\PageHashFooter{\PageMDfive{9}}
\section{Step 1 - MacroRequirement}

Un \MR{} copre una sola \textbf{responsabilita' macro durevole che il progetto deve possedere per affrontare il problema governato}. Non dice ancora come quella responsabilita' verra' realizzata, e non e' un contenitore per ogni capability importante della soluzione. Una capability non diventa MacroRequirement solo perche' e' stabile, distinguibile o tecnicamente sostituibile. Occorre verificare che il progetto debba possedere quel significato come responsabilita' macro, e non soltanto che la soluzione corrente abbia bisogno di un passaggio equivalente.

\begin{ddtaprinciple}[title={Domanda fondamentale}]
Quale responsabilita' macro deve assumersi il progetto per affrontare il problema governato? Che contributo stabile porta al problema o al boundary? Il progetto deve possedere proprio questa responsabilita', oppure essa esiste come capability distinguibile solo perche' la soluzione corrente ha scelto di decomporre il problema in questo modo? Quale parte del suo significato resta necessaria anche se cambia l'architettura con cui il problema viene affrontato?
\end{ddtaprinciple}

Il framing non e' un catalogo di MacroRequirement. Un MR deve restare \textbf{riconducibile al problema e al boundary governato}, ma non serve che compaia parola per parola nella frase di framing. Se una nuova responsabilita' cambia davvero il problema o il boundary, allora si riapre il framing.

\begin{ddtacore}[title={MR e scelte di soluzione}]
Una Decision puo' far nascere obblighi e capability reali senza per questo creare nuove responsabilita' macro. Perche' un elemento diventi MR, il suo significato deve essere governato e posseduto dal progetto al livello del problema o del boundary, indipendentemente dalla particolare decomposizione scelta per realizzarlo. La semplice possibilita' di sostituire il componente che oggi realizza una capability non e' sufficiente. Se una diversa architettura plausibile puo' affrontare lo stesso problema senza conservare quella capability come responsabilita' distinguibile, il candidato e' probabilmente una Decision, un FunctionalRequirement o un significato di livello successivo.
\end{ddtacore}

\begin{ddtaworking}[title={Confine con Decision}]
Una Decision restringe un solo MR fissando una posizione progettuale significativa: per esempio una policy, un responsibility boundary, una strategia o una scelta tecnologica/architetturale con conseguenze downstream. Se il candidato descrive gia' \emph{come} una responsabilita' viene governata o realizzata attraverso una posizione di questo tipo, va valutato come Decision e non come nuovo MR.
\end{ddtaworking}

\begin{ddtaanti}[title={Due errori opposti}]
Da una parte, non ogni sotto-problema introdotto dalla soluzione merita di diventare MR. Dall'altra, responsabilita' macro diverse non vanno compresse in un unico MR solo perche' oggi le realizza lo stesso componente o la stessa pipeline.
\end{ddtaanti}
%</PAGE-CONTENT:9>

\clearpage
%<PAGE-DESC:10>MacroRequirement - procedura umana per individuare i candidati
%<PAGE-CONTENT:10>
\PageHashFooter{\PageMDfive{10}}
\subsection{Procedura umana: dalla responsabilita' candidata al MR}

Si parte dal framing e dal boundary: quali responsabilita' macro deve possedere il progetto per affrontare quel problema? La risposta non deve anticipare una soluzione. Nel native authoring i candidati emergono dal lavoro di definizione; nella reconstruction arrivano solo dalle fonti legacy autorizzate. In questa fase una lista anche sovrabbondante va bene.

Ogni candidato riceve un \textbf{Title} e un \textbf{Intent}. Se serve, si aggiungono Context o Scope, ma solo quando aiutano davvero a capire di quale responsabilita' si sta parlando. Prima di valutarne la forma, chiedere se il progetto deve possedere quel significato \textbf{come responsabilita' macro}, oppure se il candidato descrive una capability che esiste perche' una particolare Decision o decomposizione architetturale ha introdotto quel passaggio. Nel secondo caso il significato va preservato, ma valutato a un livello inferiore. In \textbf{RECONSTRUCTION} questa ownership non deve essere dichiarata letteralmente dalla fonte: deve risultare dal significato che la fonte assegna stabilmente al progetto, senza completarlo per inferenza. I candidati si confrontano poi tra loro: overlap, containment, possibili dipendenze e diversa profondita' di decomposizione contano piu' della descrizione isolata del singolo candidato. Prima di compilare il documento si classifica: \texttt{KEEP}, \texttt{MERGE}, \texttt{SPLIT}, \texttt{LOWER LEVEL}, \texttt{REWORK} oppure \texttt{DEFER}. Solo i \texttt{KEEP} passano alla stesura completa.

\begin{ddtaworking}[title={Forma di lavoro minima per un candidato}]
\begin{lstlisting}
Candidate title:
...

Intent:
...

Context:                 [se necessario per disambiguare]
...

Scope note:              [se necessario per separare ownership]
...

Possible dependency:     [se emerge una dipendenza non gerarchica]
...
\end{lstlisting}
Questa e' una worksheet di authoring, non un nuovo tipo documentale L1 e non sostituisce il MacroRequirement finale.
\end{ddtaworking}
%</PAGE-CONTENT:10>

\clearpage
%<PAGE-DESC:11>MacroRequirement - test di separazione, split e merge
%<PAGE-CONTENT:11>
\PageHashFooter{\PageMDfive{11}}
\subsection{Test di separazione e classificazione}

Prima di accettare un candidato come MR, verificarlo con un insieme di domande complementari. Nessuna singola domanda decide da sola il livello semantico.

\begingroup\small
\begin{tabularx}{\textwidth}{L{42mm}Y}
\toprule
\textbf{Test} & \textbf{Domanda operativa} \\
\midrule
Contributo al problema & Che contributo macro porta al problema o al boundary? \\
Resilienza alla soluzione & Se venisse adottata un'architettura plausibile diversa per affrontare lo stesso problema, questa responsabilita' continuerebbe a esistere come responsabilita' macro distinguibile? Oppure esiste solo perche' l'architettura corrente ha introdotto quel particolare passaggio? \\
Singolarita' & E' una responsabilita' sola o ne sta nascondendo due che potrebbero cambiare separatamente? \\
Valore distinto & Ha un significato proprio, oppure riscrive il framing o duplica un altro MR? \\
Dependency vs containment & Il significato governato di questa responsabilita' richiede il contratto semantico di un'altra responsabilita', pur mantenendo ownership distinta? Oppure esiste soltanto un rapporto di containment, ordine operativo, dataflow o consumo di output? \\
Coerenza della futura famiglia & Le Decision che la restringerebbero appartengono allo stesso owner semantico? Il fatto che oggi non siano ancora definite non e' un problema. \\
Orizzonte di decomposizione & Le parti del candidato possono essere ristrette dallo stesso insieme di commitment governati, oppure una parte possiede gia' una Decision family supportata mentre un'altra deve fermarsi al livello MR? Una differenza persistente segnala un possibile hidden split; non giustifica il far procedere tutto il candidato insieme. \\
Stabilita' temporale & Regge oltre una singola release, configurazione o iterazione? \\
\bottomrule
\end{tabularx}

\begin{ddtacore}[title={Come leggere il controfattuale}]
Non stiamo cercando una responsabilita' che esisterebbe in ogni soluzione immaginabile. Il controfattuale deve pero' essere abbastanza profondo da non limitarsi alla sostituzione locale del componente corrente. Chiedere se ``classificare un'immagine'' sopravvive al passaggio da EfficientNet a un altro classificatore non basta: entrambe le alternative possono appartenere alla stessa decomposizione architetturale. Occorre chiedere se, adottando una diversa architettura plausibile per affrontare lo stesso problema di progetto, quel significato rimarrebbe una responsabilita' macro che il progetto deve possedere. ``Retraining EfficientNet-B4'' resta chiaramente legato alla soluzione; ``evoluzione controllata del comportamento del sistema'' resta invece riconoscibile come responsabilita' anche se cambiano tecnologia e meccanismo di adattamento. Superare il test di resilienza e' necessario ma non sufficiente per diventare MR: il candidato deve superare anche il test di ownership macro e gli altri test di separazione. Nel percorso \textbf{RECONSTRUCTION} il controfattuale serve solo a verificare la dipendenza del significato dalla decomposizione corrente; non autorizza a progettare o assumere architetture alternative non presenti nelle fonti.
\end{ddtacore}

\begin{ddtaanti}[title={Dipendenza non e' contenimento}]
Un MR puo' dipendere semanticamente da un altro MR. La presenza di una dipendenza non giustifica il merge se le due responsabilita' mantengono ownership e significato distinti. Ordine di esecuzione, passaggio di dati, consumo di un output o collocazione nella stessa pipeline non sono da soli sufficienti per stabilire \texttt{dependsOn}. La relazione va dichiarata quando il significato governato del MR dipendente fa affidamento sul contratto semantico dell'altro MR. Un indizio forte di dipendenza esiste quando il significato o le condizioni di validita' della responsabilita' dipendente non possono essere espressi correttamente senza assumere il significato governato dall'altro MR. Se la fonte mostra soltanto che un output viene consumato successivamente, senza governare la relazione semantica tra le due responsabilita', \texttt{dependsOn} non va inferito automaticamente.
\end{ddtaanti}
\endgroup
%</PAGE-CONTENT:11>

\clearpage
%<PAGE-DESC:12>MacroRequirement - campi documentali e regole di compilazione
%<PAGE-CONTENT:12>
\PageHashFooter{\PageMDfive{12}}
\subsection{Campi del MacroRequirement}

Una volta superata la separazione iniziale, il candidato viene scritto nella forma L2 completa. Ogni campo deve aggiungere significato governato; la presenza di uno slot nel template non autorizza a inventare contenuto.

\begin{tabularx}{\textwidth}{L{35mm}L{33mm}Y}
\toprule
\textbf{Campo} & \textbf{Classe} & \textbf{Uso corretto} \\
\midrule
Title & REQUIRED-CONTENT & Nome leggibile della responsabilita' macro. Non un requisito operativo, non una soluzione. \\
Intent & REQUIRED-CONTENT & Dice quale responsabilita' macro assume il progetto e quale contributo porta. \\
Context & REQUIRED-CONTENT & Il minimo background per interpretare la responsabilita' senza ambiguita'. Non e' una parafrasi dell'Intent. \\
Stakeholders & REQUIRED-SLOT & Ruoli di dominio che beneficiano della responsabilita', la governano, la operano o ne sono materialmente coinvolti. ``--'' se non emerge un significato distinto. \\
Scope & REQUIRED-SLOT & Boundary IN/OUT, quando serve a prevenire overlap o responsibility leakage. ``--'' se non aggiunge nulla. \\
Assumptions / Constraints & REQUIRED-SLOT & Solo condizioni macro valide per l'intero branch. ``--'' se non ce ne sono. \\
\texttt{dependsOn} & REQUIRED-SLOT & Dipendenza semantica non gerarchica da un altro MR: il significato della responsabilita' fa affidamento sul contratto governato dell'altro MR. Ordine operativo, dataflow o semplice consumo di output non bastano. \texttt{None} se nessuna dipendenza semantica e' governata. \\
Lifecycle / Authority & governance metadata & Stato e authority del documento, quando il corpus li usa. Non sostituiscono il significato del MR. \\
\bottomrule
\end{tabularx}

\begin{ddtacore}[title={Quando un campo non va riempito}]
Se un campo non porta niente di nuovo, non va riempito riscrivendo Intent o Context con altre parole. Dove l'assenza e' ammessa, ``--'' e' meglio di una frase di comodo.
\end{ddtacore}

\begin{ddtaanti}[title={Scope non e' un elenco di cose escluse}]
Non serve elencare modelli, pipeline o soglie solo per dire che non appartengono all'MR. Le scelte di soluzione restano semplicemente ai livelli che le governano.
\end{ddtaanti}
%</PAGE-CONTENT:12>

\clearpage
%<PAGE-DESC:13>MacroRequirement - template, split merge e STOP AT MR
%<PAGE-CONTENT:13>
\PageHashFooter{\PageMDfive{13}}
\subsection{Template L2 e risultato della review}

\begin{lstlisting}
# <MR-ID> - <Title>
Lifecycle: <...>
Authority: <...>   [se usato dalla governance del corpus]

## Intent
...
## Context
...
## Stakeholders
...
## Scope
...
## Assumptions / Constraints
...
## dependsOn
<MR-ID or None>
\end{lstlisting}

A compilazione fatta, split e merge vanno ricontrollati. Output diversi non bastano per tenere due MR separati. Condividere oggi un componente, dei dati o una sequenza operativa non basta per fonderli. Va controllato anche l'\textbf{orizzonte di decomposizione}. Se una parte del candidato dispone di commitment governati che permettono di procedere a Decision mentre un'altra parte resta semanticamente corretta solo fermandosi a MR, verificare esplicitamente se il candidato nasconde due responsabilita' distinte. Questa differenza e' un segnale di review, non una regola automatica di split: puo' anche riflettere una fonte incompleta. Un commitment governato su una parte non deve trascinare automaticamente l'altra a un livello di decomposizione che la fonte non supporta.

\begin{ddtaexample}[title={DermaTriage - STOP AT MR}]
La documentazione supporta una responsabilita' macro di indicazione della destinazione specialistica, ma non governa in modo sufficiente il vocabolario delle destinazioni, la regola di selezione, gli eventuali input specifici, il fallback o la responsabilita' di booking/assignment. La presenza, nello stesso Adaptation Layer, di commitment governati relativi alla priorita' operativa non autorizza a usarli come Decision della responsabilita' specialistica. La risposta corretta e' mantenere il significato specialistico disponibile e fermare quel branch: non si inventano Decision o FunctionalRequirement per completare la gerarchia.
\end{ddtaexample}

\begin{ddtacheck}[title={Gate D1 - MR}]
\textbf{PASS} quando la responsabilita' e' stabile, distinta dagli altri MR, collocata al livello corretto e sufficientemente chiara da non ammettere letture materialmente diverse. Dopo PASS chiedere se, nello scope corrente, esistono ulteriori commitment governati che restringono specificamente questa responsabilita' semanticamente omogenea e consentono di derivarne una Decision senza importare significato da un'altra parte del candidato, da un MR vicino o dalla struttura della soluzione.\par
\textbf{NO -> STOP AT MR.}\qquad \textbf{YES -> procedere alle Decision.}\par
Un commitment che restringe soltanto una parte di un candidato aggregato non autorizza l'intero candidato a procedere. In quel caso occorre riesaminare prima lo split/merge.
\end{ddtacheck}
%</PAGE-CONTENT:13>

\clearpage
%<PAGE-DESC:14>MacroRequirement - sufficienza semantica e diagnosi dei problemi
%<PAGE-CONTENT:14>
\PageHashFooter{\PageMDfive{14}}
\subsection{Sufficienza semantica e diagnosi}

Il gate non misura la lunghezza del testo. Chiede se, con il significato governato a disposizione, restano in piedi letture materialmente diverse della responsabilita', del suo boundary o del suo owner. In reconstruction conviene classificare le proposition critiche come:
\begin{lstlisting}
AFFIRMED
DENIED
NOT SPECIFIED
CONFLICTING
\end{lstlisting}
\texttt{NOT SPECIFIED} non e' FALSE, e \texttt{CONFLICTING} non autorizza a scegliere la lettura piu' comoda.

\begin{tabularx}{\textwidth}{L{43mm}Y}
\toprule
\textbf{Diagnosi} & \textbf{Cosa segnala} \\
\midrule
Source / documentation gap & Il significato manca, e' ambiguo o conflittuale gia' nella fonte. \\
Authoring-guide problem & Il significato si puo' rappresentare, ma domande e gate non guidano l'autore in modo stabile. \\
Metamodel pressure & Il significato e' chiaro e governato; sono i costrutti DDTA a non riuscire a rappresentarlo. \\
Representation / L2 issue & L1 basta, ma la forma documentale non e' chiara o pratica. \\
Downstream / BA issue & La documentazione regge, ma una rappresentazione analitica a valle perde significato. \\
No gap & Differenza presentazionale, oppure dettaglio fuori dallo scope governato. \\
\bottomrule
\end{tabularx}

\begin{ddtacore}[title={Il progetto non si corregge per far funzionare la metodologia}]
Un dubbio emerso in review non autorizza a cambiare documentazione, guida o metamodel. Prima va diagnosticato dove sta il problema.
\end{ddtacore}
%</PAGE-CONTENT:14>

\clearpage
%<PAGE-DESC:15>MacroRequirement - profilo LLM e prompt canonico parte 1
%<PAGE-CONTENT:15>
\PageHashFooter{\PageMDfive{15}}
\subsection{Profilo di esecuzione LLM}

Il profilo LLM segue gli stessi criteri dell'autore umano. Rende solo espliciti i controlli che servono a evitare tre errori ricorrenti: promuovere a MR un dettaglio di soluzione, chiamare decomposizione una parafrasi, riempire tutti i campi anche dove il significato non e' governato.

\subsubsection{Prompt canonico - MacroRequirement}
\begin{lstlisting}[basicstyle=\ttfamily\footnotesize]
ROLE

You are helping with DDTA documentation authoring. The DDTA Documentation
Authoring Guide is the methodological authority. Work only from the project
meaning authorized for this authoring step.

AUTHORING MODE

The caller must state one mode: NATIVE or RECONSTRUCTION.

NATIVE: start from the governed project problem, boundary and current project
meaning. Implementation state does not become project authority just because it
exists.

RECONSTRUCTION: recover project meaning only from the authorized legacy sources.
Missing meaning is not repaired with domain knowledge, earlier DDTA
reconstructions, Base Analysis, threat models, tooling needs or implementation
assumptions.

TASK

Identify, separate, and author MacroRequirements. A MacroRequirement owns one
durable macro responsibility the project must possess to address the governed
problem within its boundary. It is not a container for every stable capability
or important project concern.
\end{lstlisting}
%</PAGE-CONTENT:15>

\clearpage
%<PAGE-DESC:16>MacroRequirement - prompt canonico LLM parte 2
%<PAGE-CONTENT:16>
\PageHashFooter{\PageMDfive{16}}
\begin{lstlisting}[basicstyle=\ttfamily\scriptsize]
CANDIDATE DISCOVERY

1. Identify candidate responsibilities without assuming that every stable
   capability is a MacroRequirement.
2. Give each candidate at least a Title and Intent before judging it.
3. Add Context, Scope, or a possible dependency only when they are needed to
   understand ownership or separation.
4. Ask whether the project must own that meaning as a macro responsibility, or
   whether it exists because the current solution decomposes another
   responsibility into that capability.
5. Preserve solution-induced obligations, but classify them LOWER_LEVEL when a
   technical or architectural commitment introduced that decomposition.

SEPARATION TESTS

For every candidate ask:
- What distinct macro contribution does it make to the governed problem or
  boundary?
- MACRO OWNERSHIP: must the project own this meaning as a macro responsibility,
  or does it exist because the current solution decomposes another
  responsibility into this capability? In RECONSTRUCTION, infer this only from
  source-supported project meaning; do not require literal ownership wording.
- SOLUTION RESILIENCE: if a plausibly different architecture were used to
  address the same project problem, would this meaning still exist as a
  distinguishable macro responsibility? Do not test only replacement of the
  current component by an equivalent component. In RECONSTRUCTION, use this
  only as a classification counterfactual; do not design or assume unsupported
  alternative architectures.
- Is it one responsibility, or several responsibilities that could change apart?
- Does it add its own meaning, rather than restating the framing or another MR?
- Is its relation to another MR a semantic dependency, containment, operational
  sequence, dataflow, or simple output consumption?
- Would the Decisions restricting it belong to one coherent semantic owner?
- DECOMPOSITION HORIZON: can the whole candidate be restricted by the same
  governed commitment family, or can one part proceed to Decisions while
  another must STOP AT MR? Treat a mismatch as SPLIT-PRESSURE, not an automatic
  split.
- Does the responsibility remain stable beyond one release or configuration?

No single PASS makes a candidate a MacroRequirement. Solution resilience or a
coherent Decision family is not sufficient without macro ownership.

Do NOT require complete independence from other MR. A valid MR may depend on
another MR through dependsOn while retaining separate ownership. Do not infer
dependsOn from execution order, dataflow, or output consumption alone.

CLASSIFY BEFORE FULL AUTHORING

KEEP | MERGE | SPLIT | LOWER_LEVEL | REWORK | DEFER

Only KEEP candidates may become final MacroRequirements.
\end{lstlisting}

\begin{ddtaanti}[title={Regola per l'LLM}]
Conoscere un dettaglio della soluzione non gli assegna automaticamente un posto nel livello MR. Il modello deve prima stabilire quale responsabilita' governa quel significato, se il progetto deve possederlo come responsabilita' macro e se regge a un cambio plausibile della decomposizione architetturale. Un singolo test superato non basta per promuovere il candidato.
\end{ddtaanti}
%</PAGE-CONTENT:16>

\clearpage
%<PAGE-DESC:17>MacroRequirement - prompt LLM campi e schema di analisi
%<PAGE-CONTENT:17>
\PageHashFooter{\PageMDfive{17}}
\begin{lstlisting}[basicstyle=\ttfamily\scriptsize]
FULL MR AUTHORING

For each KEEP candidate produce:
- Title: REQUIRED-CONTENT
- Intent: REQUIRED-CONTENT
- Context: REQUIRED-CONTENT
- Stakeholders: REQUIRED-SLOT; use "-" if no distinct governed meaning exists
- Scope: REQUIRED-SLOT; use "-" if it adds no distinct boundary information
- Assumptions / Constraints: REQUIRED-SLOT; use "-" if none are governed
- dependsOn: REQUIRED-SLOT; use "None" when no semantic dependency is governed

FIELD NOVELTY RULE

Do not fill a field by paraphrasing another field. Scope is not a list of
technical choices to exclude. Assumptions/Constraints must apply to the whole
MR branch.

dependsOn is semantic and non-hierarchical. Do not infer it from narrative
order, pipeline order, dataflow, shared data, or simple consumption of another
MR's output. Use dependsOn only when the governed meaning of one MR relies on
the semantic contract owned by another MR. A strong indicator exists when the
dependent responsibility cannot be expressed correctly without assuming that
governed meaning.

SEMANTIC SUFFICIENCY

Ask whether materially different interpretations of the responsibility can
still survive. If yes, REWORK or STOP rather than inventing lower-level meaning.
A stable MR may PASS even when lower-level Decisions or FR are NOT SPECIFIED.

BRANCH STOP RULE

After an MR passes, proceed to Decisions only when additional governed
commitments specifically restrict that semantic responsibility.

A commitment that restricts another MR, or only one separable part of an
aggregated candidate, does not authorize this branch to proceed.

If different parts of the candidate have different decomposition horizons,
re-open SPLIT/MERGE before deciding PROCEED TO DECISIONS. A mismatch is a
review signal, not an automatic split.

Otherwise STOP AT MR.
\end{lstlisting}

\subsection{Worksheet richiesta all'LLM}
La worksheet rende verificabile la classificazione dei candidati e resta separata dalla documentazione governata finale.
%</PAGE-CONTENT:17>

\clearpage
%<PAGE-DESC:18>MacroRequirement - formato esatto output LLM e gate finale
%<PAGE-CONTENT:18>
\PageHashFooter{\PageMDfive{18}}
\subsubsection{Formato esatto richiesto}
\begin{lstlisting}[basicstyle=\ttfamily\scriptsize]
[DDTA MR ANALYSIS]
AUTHORING MODE: NATIVE | RECONSTRUCTION
AUTHORIZED PROJECT BASIS:
- ...

CANDIDATE <n>
TITLE: ...
INTENT: ...
SOURCE / GOVERNED SUPPORT: ...
CONTEXT NOTE: ... | -
SCOPE NOTE: ... | -
POSSIBLE DEPENDENCY: ... | None
MACRO-OWNERSHIP: PASS | LOWER_LEVEL | UNCERTAIN
SOLUTION-RESILIENCE: PASS | REWORK
DISTINCT-RESPONSIBILITY: PASS | MERGE | SPLIT | REWORK
DECOMPOSITION-HORIZON: COHERENT | SPLIT-PRESSURE | NOT-YET-TESTABLE
SEMANTIC-LEVEL: MR | LOWER_LEVEL | UNCERTAIN
CLASSIFICATION: KEEP | MERGE | SPLIT | LOWER_LEVEL | REWORK | DEFER
RATIONALE: ...

[FINAL MACROREQUIREMENTS]
MR <id or candidate-id> - <Title>
Intent: ...
Context: ...
Stakeholders: ... | -
Scope: ... | -
Assumptions / Constraints: ... | -
dependsOn: <MR-ID> | None
MR VERDICT: PASS | REWORK
NEXT: PROCEED TO DECISIONS | STOP AT MR
\end{lstlisting}

\begin{ddtacheck}[title={Controlli prima di accettare l'output LLM}]
Prima di accettare l'output:
\begin{itemize}
  \item ogni MR finale viene da un candidato classificato \texttt{KEEP};
  \item ogni \texttt{KEEP} ha superato il macro-ownership test; solution resilience da sola non basta;
  \item nessun dettaglio marcato \texttt{LOWER\_LEVEL} rientra come responsabilita' macro senza una nuova giustificazione;
  \item il solution-resilience test considera un cambio plausibile di architettura, senza progettare alternative non autorizzate in reconstruction;
  \item split e merge sono motivati sul significato, non sulla struttura attuale del software;
  \item una differenza di decomposition horizon viene riesaminata come possibile split, mai trattata come split automatico;
  \item \texttt{dependsOn} non viene dedotto dal solo ordine operativo, dataflow o consumo di output;
  \item Context e Scope portano informazione nuova invece di ripetere l'Intent;
  \item \texttt{STOP AT MR} si applica localmente alla responsabilita': commitment presenti in un'altra parte del candidato non autorizzano il branch a procedere;
  \item l'output LLM passa comunque dalla stessa review semantica dell'autore umano.
\end{itemize}
\end{ddtacheck}
%</PAGE-CONTENT:18>

\clearpage
%<PAGE-DESC:19>Decision - definizione confine e discovery
%<PAGE-CONTENT:19>
\PageHashFooter{\PageMDfive{19}}
\section{Step 2 - Decision}
\subsection{Che cos'è una Decision}
Una \textbf{Decision} è un commitment progettuale significativo che restringe esattamente un MacroRequirement, attraverso una posizione, una policy, un boundary, una convenzione, una strategia, un placement o una scelta architetturale governata. Non sostituisce la responsabilità macro dell'MR e non anticipa gli obblighi operativi che apparterranno ai FunctionalRequirement.

\begin{ddtaprinciple}[title={Confine MR - Decision}]
\textbf{MR}: quale responsabilità macro possiede il progetto?\par
\textbf{Decision}: quale scelta governata restringe quella responsabilità senza cambiarne l'identità?
\end{ddtaprinciple}

\subsection{Discovery e promotion sono due passaggi diversi}
Prima si cercano i significati che potrebbero essere Decision; solo dopo li si sottopone ai test di promozione. Un boundary interessante, una tecnologia importante o un comportamento ripetuto non diventano Decision solo perché sono stati scoperti.

\begin{ddtacheck}[title={Domande per la discovery}]
\begin{itemize}
\item La fonte contiene una policy, una convenzione, una strategia o una scelta architetturale?
\item Esiste una separazione di fasi, un'integrazione, un lifecycle o un placement con significato proprio?
\item Una condizione è accompagnata da una risposta progettuale selezionata?
\item Un dettaglio tecnico nasconde un significato più astratto che resta comprensibile anche senza il nome concreto?
\end{itemize}
\end{ddtacheck}
%</PAGE-CONTENT:19>

\clearpage
%<PAGE-DESC:20>Decision - worksheet candidate e source support
%<PAGE-CONTENT:20>
\PageHashFooter{\PageMDfive{20}}
\subsection{Worksheet minima della candidate}
\begin{lstlisting}
Candidate ID: ...
Parent MR: ...
Source evidence: ...
Minimum source-supported meaning: ...
Possible Decision meaning: ...
Possible lower-level owner: ...
Open source gaps: ...
\end{lstlisting}

\begin{ddtacore}[title={Source-first}]
La candidate parte dal significato minimo realmente sostenuto dalla documentazione autorizzata. Il metodo distingue sempre ciò che la fonte stabilisce, ciò che non stabilisce e ciò che resta ambiguo.
\end{ddtacore}

\begin{ddtaexample}[title={DermaTriage - fallback senza immagine}]
La documentazione descrive una condizione senza immagine e un percorso basato sui sintomi. La candidate non nasce dalla parola ``fallback'', ma dalla scelta progettuale più astratta: continuare la valutazione di triage anche quando l'immagine non è disponibile.
\end{ddtaexample}

\begin{ddtacheck}[title={Domande da annotare, non da risolvere per intuizione}]
Quali output sono garantiti nel percorso senza immagine? Quali informazioni sintomatologiche sono necessarie? La fonte richiede equivalenza con il percorso image-based? Se la documentazione non risponde, il dubbio resta visibile.
\end{ddtacheck}
%</PAGE-CONTENT:20>

\clearpage
%<PAGE-DESC:21>Decision - MR stability neutralization semantic level
%<PAGE-CONTENT:21>
\PageHashFooter{\PageMDfive{21}}
\subsection{Tre test prima della governance}
\subsubsection{MR-stability}
Chiedere se il significato può cambiare mantenendo stabile la responsabilità macro dell'MR. Il confronto va fatto con l'identità semantica dell'MR, non con il vocabolario o la realizzazione corrente.

\subsubsection{Neutralization}
Rimuovere temporaneamente nomi di prodotti, modelli, endpoint, file, protocolli, parametri, soglie e binding concreti. Se non resta alcun significato progettuale, la candidate è normalmente lower-level.

\begin{ddtaanti}[title={Neutralization non è promotion}]
Se dopo la neutralizzazione resta una strategia o un boundary comprensibile, significa solo che esiste una candidate significativa. Non basta ancora a dimostrare che il progetto la governi come Decision autonoma.
\end{ddtaanti}

\subsubsection{Semantic-level test}
Confrontare la candidate con MR identity, FunctionalRequirement, realization, interface detail, configuration, parameter/binding, data/model fact, verification evidence e source gap.

\begin{ddtaexample}[title={DermaTriage - retrieval storico}]
\texttt{ChromaDB + embedding model + cosine + top-5} può neutralizzarsi in ``historical-case retrieval''. La candidate sopravvive; resta però da capire se quel significato possiede governance autonoma o resta comportamento interno della pipeline.
\end{ddtaexample}
%</PAGE-CONTENT:21>

\clearpage
%<PAGE-DESC:22>Decision - independent governance e removal replacement
%<PAGE-CONTENT:22>
\PageHashFooter{\PageMDfive{22}}
\subsection{Independent-governance test}
La domanda centrale è: \emph{quale evidenza nella documentazione mostra che il progetto possiede questa scelta proprio a questo livello semantico, invece di documentare soltanto la realizzazione corrente?}

Evidenza favorevole può venire da framing architetturale o di policy, rationale, alternative deliberate, boundary autonomi, conseguenze proprie, stabilità cross-document e significato di change/review. La ripetizione della stessa realizzazione in più punti non basta da sola.

\subsection{Removal / replacement significance}
Chiedere: se il significato neutralizzato venisse rimosso o sostituito lasciando stabile l'MR, la documentazione supporterebbe la conclusione che è cambiato un commitment governato, oppure soltanto la sua realizzazione?

\begin{ddtaexample}[title={Esempi di contrasto}]
Cambiare il nome di un endpoint non equivale a rimuovere una strategia di accesso. Cambiare un modello concreto non equivale a rimuovere uno stadio architetturale governato. Cambiare un parametro non equivale a cambiare una policy.
\end{ddtaexample}

\begin{ddtacheck}[title={Domanda utile nella review}]
Se la fonte documenta due percorsi di accesso, stabilisce anche che entrambi debbano essere preservati come scelta di progetto, oppure descrive soltanto l'interfaccia corrente?
\end{ddtacheck}
%</PAGE-CONTENT:22>

\clearpage
%<PAGE-DESC:23>Decision - decomposition, semantic owner e utilita' downstream
%<PAGE-CONTENT:23>
\PageHashFooter{\PageMDfive{23}}
\subsection{Decomposition e semantic owner}
Una candidate non deve assorbire significati che possono cambiare o essere valutati indipendentemente. Valutare lo split quando una sub-parte:
\begin{itemize}
\item può variare indipendentemente;
\item ha un semantic owner diverso;
\item può essere revisionata o verificata separatamente;
\item può cambiare senza cambiare il significato parent.
\end{itemize}
Questi segnali aiutano a individuare significati distinti, ma \textbf{non impongono da soli un branch Decision separato}. La struttura documentale deve restare utile anche dopo la derivazione dei requisiti.

\begin{ddtaprinciple}[title={Semantic owner non equivale automaticamente a branch autonomo}]
Una differenza concettuale può essere reale senza meritare una Decision autonoma. Dopo l'authoring degli FR, ogni branch deve dimostrare una conseguenza downstream distinguibile per implementazione, verifica, modifica indipendente o analisi. Se il significato vive soltanto come parte del contratto operativo di un branch sibling, riesaminare \texttt{MERGE} o \texttt{REWORK}; non creare un FR artificiale per giustificare lo split.
\end{ddtaprinciple}

\begin{ddtaexample}[title={DermaTriage - P-scale}]
Il dominio P1-P4 e la regola urgency-to-priority sono significati distinguibili: il primo definisce il codominio della priorita', la seconda definisce come scegliere il livello. Questo non prova pero' che richiedano due branch Decision. Se il dominio non produce alcuna conseguenza downstream autonoma e diventa operativo soltanto nel branch che governa la derivazione, il significato va preservato in quel contratto senza mantenere una diramazione documentale separata.
\end{ddtaexample}

\subsection{Campi della Decision}
\begin{lstlisting}
Parent MR: <exactly one MR>
Context: <situazione progettuale o di dominio necessaria a capire perche' la scelta e' rilevante>
Decision: <the governed commitment>
Consequences: <what follows from adopting the choice>
\end{lstlisting}

\begin{ddtacheck}[title={Regola per il Context}]
Il Context descrive il progetto, il dominio o la condizione in cui la scelta si applica. Non descrive il processo di analisi, non cita la documentazione come soggetto e non giustifica perché l'analista ha classificato l'elemento come Decision.
\end{ddtacheck}
%</PAGE-CONTENT:23>

\clearpage
%<PAGE-DESC:24>Decision - esempio completo e MR prose leakage
%<PAGE-CONTENT:24>
\PageHashFooter{\PageMDfive{24}}
\subsection{Esempio completo}
\begin{ddtaexample}[title={DermaTriage - continuità del triage senza immagine}]
\textbf{Parent MR}: MR-01.\par
\textbf{Context}: un caso dermatologico può avere un'immagine della lesione, oppure no.\par
\textbf{Decision}: DermaTriage valuta il triage anche se manca l'immagine della lesione. In quel caso usa le informazioni sui sintomi disponibili sul caso.\par
\textbf{Consequences}: se l'immagine manca, il triage si può comunque fare. Il percorso basato sui sintomi può usare informazioni diverse; non deve per forza dare gli stessi risultati del percorso con immagine.
\end{ddtaexample}

\subsection{MR-prose Decision leakage}
Durante la review, cercare strutture del tipo \texttt{condizione + risposta progettuale selezionata} dentro Context o Scope dell'MR. Se la risposta può cambiare mentre la responsabilità macro resta stabile, il testo MR può contenere una Decision che va spostata al proprio semantic owner.

\begin{ddtacheck}[title={Domanda di controllo}]
La frase nel Context descrive davvero il contesto della responsabilità, oppure prescrive già come il progetto risponde a una condizione particolare?
\end{ddtacheck}
%</PAGE-CONTENT:24>

\clearpage
%<PAGE-DESC:25>Decision - family regression, downstream utility e closure
%<PAGE-CONTENT:25>
\PageHashFooter{\PageMDfive{25}}
\subsection{Decision-family regression}
Dopo la review individuale, controllare la famiglia nel suo insieme:
\begin{itemize}
\item pairwise non-overlap e frammentazione inutile;
\item candidate troppo ampie o troppo tecniche;
\item significati governati rimasti senza owner;
\item boundary di fase, integrazione, lifecycle e data relationship;
\item MR-prose leakage;
\item completezza rispetto alla source corrente.
\end{itemize}
Questa prima regression può chiudere la discovery delle Decision, ma la \textbf{struttura dei branch resta provvisoria fino alla derivazione degli FR}.

\begin{ddtacheck}[title={Downstream Utility Test - regression obbligatoria post-FR}]
Per ogni Decision che si vuole mantenere come branch autonomo, dopo la derivazione degli FR chiedere:
\begin{quote}
Se elimino questo branch come unita' autonoma e trasferisco il suo significato nel branch che possiede i comportamenti che lo utilizzano, perdo una distinzione necessaria per implementazione, verifica, modifica indipendente o analisi?
\end{quote}
\textbf{YES}: il branch possiede utilita' downstream distinta. \textbf{NO}: la differenza e' soltanto concettuale/documentale; riaprire la struttura e valutare \texttt{MERGE} o \texttt{REWORK}. Una Decision semanticamente vera non e' sufficiente, da sola, a giustificare un branch autonomo.
\end{ddtacheck}

Una candidate può quindi essere mantenuta, fusa, separata, riclassificata a un livello inferiore, rielaborata o lasciata in HOLD.

\subsection{Closure layer-specific}
\begin{lstlisting}[basicstyle=\ttfamily\small]
SOURCE INSPECTION: COMPLETE | INCOMPLETE
CANDIDATE DISCOVERY: EXHAUSTED FOR CURRENT SOURCE SET | OPEN
DECISION FAMILY: PROVISIONAL | RESOLVED | HOLD
LOWER-LEVEL SEMANTICS: OPEN | RESOLVED
FR AUTHORING: NOT STARTED | STARTED | COMPLETE
POST-FR DOWNSTREAM UTILITY: NOT RUN | PASS | REWORK
\end{lstlisting}

\begin{ddtacheck}[title={Domanda di chiusura}]
La family e' \texttt{RESOLVED} solo quando nessun dubbio cambia identita' o numero delle Decision e, se gli FR sono completi, il Downstream Utility Test non lascia branch autonomi privi di utilita' distinta.
\end{ddtacheck}
%</PAGE-CONTENT:25>

\clearpage
%<PAGE-DESC:26>Decision - LLM execution profile prompt canonico parte 1
%<PAGE-CONTENT:26>
\PageHashFooter{\PageMDfive{26}}
\subsection{Profilo di esecuzione LLM}
Il modello applica la stessa procedura dell'autore umano. Non crea authority, non completa source gap per plausibilita' e non trasforma una candidate in Decision perche' ``sembra architetturale''.

\subsubsection{Prompt canonico - Decision}
\begin{lstlisting}[basicstyle=\ttfamily\scriptsize]
ROLE
You are helping with DDTA documentation authoring. Work on exactly one accepted
MacroRequirement at a time. The authorized project basis is the only authority
for project facts and commitments. The LLM is assistive and cannot create
missing project meaning.

AUTHORING MODE
The caller must state NATIVE or RECONSTRUCTION.
In RECONSTRUCTION, recover meaning only from authorized legacy sources. Do not
repair gaps using domain knowledge, previous DDTA answers, Base Analysis,
threat models, tooling needs, or plausible design.

TASK
Discover, classify, and author the Decision family that restricts the supplied
parent MR. Do not author FunctionalRequirements.

SOURCE COVERAGE
List the authorized sources inspected. Do not claim discovery completeness
until the declared source set has been checked.

CANDIDATE DISCOVERY
Search for source-supported policy, convention, strategy, architecture or
placement boundaries, phase or integration boundaries, lifecycle/data
relationships, condition+selected-response structures, and technical choices
that may expose a more abstract governed meaning. Discovery is not promotion.
\end{lstlisting}
%</PAGE-CONTENT:26>

\clearpage
%<PAGE-DESC:27>Decision - prompt LLM candidate gates
%<PAGE-CONTENT:27>
\PageHashFooter{\PageMDfive{27}}
\begin{lstlisting}[basicstyle=\ttfamily\scriptsize]
FOR EVERY CANDIDATE
A. SOURCE-SUPPORTED MEANING
Separate ESTABLISHED / NOT ESTABLISHED / AMBIGUITY / SOURCE GAP.

B. MR STABILITY
Could this meaning change while the parent MR remains the same macro
responsibility? PASS | FAIL | HOLD

C. NEUTRALIZATION
Remove concrete product/model/endpoint/protocol/file names, parameter values,
thresholds and configuration bindings. Does meaningful project strategy,
boundary, policy, convention or placement remain? PASS | FAIL | N/A | HOLD
Surviving neutralization creates a candidate; it does not prove Decision status.

D. SEMANTIC LEVEL
Compare against MR identity, FR/operational obligation, realization, interface,
configuration, parameter/binding, data/model fact, verification evidence,
observed result and source gap. PASS | FAIL | HOLD

E. INDEPENDENT GOVERNANCE
What original-source evidence establishes that the surviving meaning is owned
and governed by the project at this semantic level rather than merely repeated
as current realization? PASS | FAIL | HOLD

F. REMOVAL / REPLACEMENT SIGNIFICANCE
If the meaning were removed or replaced while the MR stayed stable, would the
authorized basis support that a governed project commitment changed rather
than only its realization? PASS | FAIL | HOLD
\end{lstlisting}
%</PAGE-CONTENT:27>

\clearpage
%<PAGE-DESC:28>Decision - prompt LLM decomposition, downstream regression e output
%<PAGE-CONTENT:28>
\PageHashFooter{\PageMDfive{28}}
\begin{lstlisting}[basicstyle=\ttfamily\scriptsize]
DECOMPOSITION / SEMANTIC OWNER
Identify separable meanings, but do not create a separate Decision branch only
because a meaning has a different conceptual owner. Branch autonomy must later
survive the post-FR downstream-utility regression.

PROVISIONAL CLASSIFICATION
KEEP | MERGE | SPLIT | LOWER_LEVEL | REWORK | BOUNDARY_REVIEW |
HOLD / SOURCE_INSUFFICIENT

[FAMILY REGRESSION]
PAIRWISE NON-OVERLAP: PASS | REWORK
UNOWNED GOVERNED MEANING: NONE | ...
BOUNDARY FINDINGS: ...
MR-PROSE LEAKAGE: NONE | ...
FAMILY ACTIONS: ...
POST-FR DOWNSTREAM UTILITY: NOT_TESTABLE_YET | PASS | REWORK

[POST-FR DOWNSTREAM UTILITY TEST]
For each autonomous Decision branch ask:
"If this branch is removed as an autonomous unit and its governed meaning is
 transferred to the sibling branch that owns the behaviors using it, is any
 distinction required for implementation, verification, independent change,
 or analysis lost?"
YES -> autonomous branch may remain.
NO  -> reopen MERGE / REWORK. Do not invent an FR to save the branch.

[CLOSURE]
SOURCE INSPECTION: COMPLETE | INCOMPLETE
CANDIDATE DISCOVERY: EXHAUSTED FOR CURRENT SOURCE SET | OPEN
DECISION FAMILY: PROVISIONAL | RESOLVED | HOLD
LOWER-LEVEL SEMANTICS: OPEN | RESOLVED
FR AUTHORING: NOT STARTED | STARTED | COMPLETE

[FINAL DECISIONS]
DEC-<id> - <title>
Parent MR: ...
Context: ...
Decision: ...
Consequences: ...
\end{lstlisting}

\begin{ddtaanti}[title={Stop rule}]
Do not author FunctionalRequirements during Decision discovery. Do not complete source gaps. Do not preserve a Decision branch merely because its statement is true or conceptually distinct. Final branch autonomy requires the post-FR Downstream Utility Test and human semantic review.
\end{ddtaanti}
%</PAGE-CONTENT:28>

\clearpage
%<PAGE-DESC:29>FunctionalRequirement - definizione, termini e domanda fondamentale
%<PAGE-CONTENT:29>
\PageHashFooter{\PageMDfive{29}}
\section{Step 3 - FunctionalRequirement}

\subsection{Che cos\textquotesingle è un
FunctionalRequirement}

Un FunctionalRequirement (FR) descrive un comportamento operativo che il
progetto richiede affinché una Decision possa essere realmente
rispettata.

La Decision stabilisce una scelta, una regola o un vincolo progettuale.
Il FunctionalRequirement chiarisce che cosa il sistema deve fare,
produrre o garantire sotto quella scelta.

Per questo un FR non dovrebbe limitarsi a riscrivere la Decision con
parole diverse. Allo stesso tempo, non dovrebbe essere reso così
astratto da perdere il comportamento che il progetto ha realmente
documentato.

\textbf{Domanda fondamentale}

\begin{ddtacheck}[title={Domanda di review}]
Quale comportamento operativo deve essere vero perché questa Decision
sia soddisfatta?\\
Quale parte di questo comportamento è effettivamente sostenuta dalle
fonti autorizzate?
\end{ddtacheck}

In questa sezione ricorrono alcuni termini DDTA:

\begin{itemize}
\tightlist
\item
  \textbf{owner del significato}: l\textquotesingle elemento documentale
  che ha la responsabilità di governare quel significato;
\item
  \textbf{authority}: la fonte o il livello di governance che ci
  autorizza a trattare un\textquotesingle informazione come significato
  del progetto;
\item
  \textbf{binding}: un collegamento concreto fra significati che devono
  rimanere associati, per esempio fra un risultato e la richiesta o il
  caso a cui appartiene;
\item
  \textbf{candidate}: un elemento ancora in valutazione, non ancora
  accettato;
\item
  \textbf{review}: il controllo umano con cui decidiamo se un candidato
  è corretto, va modificato oppure deve rimanere aperto.
\end{itemize}

L\textquotesingle obiettivo non è sostituire questi termini, ma fare in
modo che il lettore sappia sempre che cosa significano quando vengono
usati.
%</PAGE-CONTENT:29>

\clearpage
%<PAGE-DESC:30>FunctionalRequirement - confine Decision-FR, parent unico e classificazioni di review
%<PAGE-CONTENT:30>
\PageHashFooter{\PageMDfive{30}}
\subsection{Dove finisce la Decision e dove inizia il
FunctionalRequirement}

La distinzione può essere letta in modo semplice:

\begin{itemize}
\tightlist
\item
  la \textbf{Decision} risponde a: \emph{quale scelta, regola o vincolo
  progettuale stiamo adottando?}
\item
  il \textbf{FunctionalRequirement} risponde a: \emph{quale
  comportamento deve essere garantito perché quella scelta sia
  rispettata?}
\end{itemize}

Se il testo di un candidato FR può essere sostituito integralmente con
quello della Decision senza perdere alcuna informazione, probabilmente
stiamo descrivendo due volte lo stesso significato.

In quel caso non dobbiamo forzare la presenza di un FR: dobbiamo
riesaminare il confine fra i due livelli.

Se invece le fonti descrivono un comportamento operativo che resta
distinguibile dalla scelta espressa nella Decision, quel comportamento
può avere una propria identità come FunctionalRequirement.

\textbf{Domanda di controllo}

\begin{ddtacheck}[title={Domanda di review}]
Se mantengo invariata la Decision ma elimino questo FR, quale
comportamento richiesto dal progetto smette di essere governato?
\end{ddtacheck}

Se la risposta è "nessuno", il candidato FR deve essere riesaminato.

\subsection{Ogni FR ha una sola Decision
parent}

Ogni FunctionalRequirement appartiene semanticamente a una sola
Decision.

Il MacroRequirement di riferimento si ricava risalendo dalla Decision.
Non costituisce quindi un secondo parent del FR.

Durante la ricerca delle Decision possono però emergere significati
reali che, dopo la review, non vengono accettati come Decision autonome.
In questo caso possiamo usare alcune classificazioni provvisorie.

\textbf{Riclassificato a un livello più basso
(\passthrough{\lstinline!LOWER\_LEVEL!})}

Il significato è sostenuto dalle fonti, ma non possiede sufficiente
autonomia per essere una Decision. L\textquotesingle informazione non
viene eliminata: deve essere conservata e collocata nel livello o presso
l\textquotesingle owner più appropriato.

\textbf{Lasciato in sospeso (\passthrough{\lstinline!HOLD!})}

Il significato non viene né accettato né scartato. Mancano ancora
informazioni, non è chiaro il confine con altri elementi oppure non
sappiamo ancora chi debba governarlo. Il candidato resta quindi aperto
fino a una review successiva.

\passthrough{\lstinline!LOWER\_LEVEL!} e \passthrough{\lstinline!HOLD!}
\textbf{non sono tipi di Decision}. Sono esiti della review di un
candidato.

Di conseguenza, un candidato classificato
\passthrough{\lstinline!LOWER\_LEVEL!} o \passthrough{\lstinline!HOLD!}
non può diventare parent di un FR finché non viene stabilita una vera
Decision che possieda quel comportamento.
%</PAGE-CONTENT:30>

\clearpage
%<PAGE-DESC:31>FunctionalRequirement - individuazione source-first dei comportamenti
%<PAGE-CONTENT:31>
\PageHashFooter{\PageMDfive{31}}
\subsection{Come individuare gli FR partendo dalle
fonti}

Per ogni Decision già accettata conviene procedere senza cercare subito
di "scrivere requisiti".

Prima ricostruiamo il comportamento che le fonti effettivamente
governano.

\begin{enumerate}
\def\labelenumi{\arabic{enumi}.}
\tightlist
\item
  Rileggi la Decision e le fonti autorizzate che la riguardano.
\item
  Individua le affermazioni minime che descrivono comportamenti,
  risultati, condizioni o relazioni necessarie.
\item
  Chiediti quali di queste informazioni descrivono \textbf{che cosa deve
  fare il sistema} e quali descrivono invece \textbf{come il sistema è
  realizzato oggi}.
\item
  Formula uno o più candidati FR soltanto dopo aver chiarito questa
  distinzione.
\item
  Se manca un\textquotesingle informazione, registrala come domanda
  aperta invece di completarla con una supposizione plausibile.
\item
  Solo dopo valuta se i candidati devono essere mantenuti, uniti,
  separati, riclassificati o riesaminati insieme.
\end{enumerate}

Un dettaglio tecnico presente in una tabella, in un diagramma, nel
codice o in una configurazione è comunque informazione del progetto.
\textbf{Non deve essere perso.}

La sua presenza, però, non ci dice automaticamente dove debba essere
governato.

Prima di inserirlo nelle clausole normative di un FR dobbiamo capire:

\begin{itemize}
\tightlist
\item
  che cosa significa;
\item
  a quale comportamento è collegato;
\item
  chi ne possiede il significato;
\item
  quale fonte ci autorizza a trattarlo come obbligazione, realizzazione,
  configurazione o altro elemento documentale.
\end{itemize}

Questa distinzione serve a collocare correttamente
l\textquotesingle informazione, non a ridurla.
%</PAGE-CONTENT:31>

\clearpage
%<PAGE-DESC:32>FunctionalRequirement - scheda candidate, branch utility e obbligazione coerente
%<PAGE-CONTENT:32>
\PageHashFooter{\PageMDfive{32}}
\subsection{Scheda minima per analizzare un candidato
FR}

Durante la review può essere utile compilare una scheda molto semplice:

\begin{lstlisting}[basicstyle=\ttfamily\small]
ID del candidato:
Decision parent:
Evidenza originale:
Comportamento operativo candidato:
Che cosa non sarebbe piu governato se lo eliminassi:
Dettagli di realizzazione/configurazione/binding collegati:
Possibile altro owner del significato:
Informazioni mancanti o ambigue:
Esito provvisorio:
\end{lstlisting}

L\textquotesingle esito provvisorio può essere, per esempio:
\begin{itemize}
\tightlist
\item mantenere il candidato (\passthrough{\lstinline!KEEP!});
\item unirlo a un altro (\passthrough{\lstinline!MERGE!});
\item dividerlo (\passthrough{\lstinline!SPLIT!});
\item riclassificarlo a un livello più basso (\passthrough{\lstinline!LOWER\_LEVEL!});
\item rielaborarlo (\passthrough{\lstinline!REWORK!});
\item lasciarlo temporaneamente in sospeso (\passthrough{\lstinline!HOLD!}).
\end{itemize}
I codici rendono ripetibile la review ma non sostituiscono la motivazione.

\begin{ddtacore}[title={Non inventare un FR per salvare una Decision}]
La gerarchia non va riempita artificialmente. Se non emerge un comportamento operativo distinto, non creare un FR solo per dare un figlio alla Decision. Al contrario, \textbf{riapri la Decision-family review}. Una Decision senza FR propri può restare autonoma solo se il Downstream Utility Test mostra una distinzione necessaria per implementazione, verifica, modifica indipendente o analisi. In assenza di tale utilita', valutare \texttt{MERGE} o \texttt{REWORK} con il branch che possiede il comportamento. Le lacune di fonte restano esplicite e non vanno colmate per supposizione.
\end{ddtacore}

\subsection{Un FR rappresenta una sola obbligazione
coerente}
La regola di base è:
\begin{lstlisting}[basicstyle=\ttfamily\small]
1 Requirement != 1 frase
1 Requirement  = 1 obbligazione normativa coerente
\end{lstlisting}

Un FR può contenere più clausole, condizioni o un output composto senza essere automaticamente diviso. La domanda non è ``quante frasi ci sono?'', ma:

\begin{ddtacheck}[title={Domanda di review}]
queste parti descrivono insieme un unico comportamento richiesto oppure rappresentano obbligazioni che il progetto può governare separatamente?
\end{ddtacheck}

Il numero di punti elenco, campi JSON, tecnologie citate, test o righe di una tabella non decide lo split.
%</PAGE-CONTENT:32>

\clearpage
%<PAGE-DESC:33>FunctionalRequirement - split e verificabilita' senza frammentazione
%<PAGE-CONTENT:33>
\PageHashFooter{\PageMDfive{33}}
\subsection{Quando separare un FR}

Valuta la creazione di due FR distinti quando una parte del candidato
possiede una vita realmente indipendente dalle altre.

Per esempio, chiediti se quella parte può:

\begin{itemize}
\tightlist
\item
  essere introdotta o rimossa senza cambiare l\textquotesingle altra;
\item
  essere modificata senza cambiare l\textquotesingle identità
  dell\textquotesingle altra;
\item
  fallire mentre l\textquotesingle altra continua a essere soddisfatta;
\item
  essere verificata e governata separatamente;
\item
  avere un owner differente;
\item
  avere un proprio ciclo di cambiamento o una propria autorità di
  review.
\end{itemize}

Nessuno di questi segnali, preso isolatamente, impone automaticamente
uno split. Servono a capire se esistono davvero due obbligazioni
autonome.

\textbf{Domanda split/merge}

\begin{ddtacheck}[title={Domanda di review}]
Se separo questa parte, ottengo un\textquotesingle obbligazione che il
progetto governa autonomamente oppure sto soltanto frammentando la
descrizione di un unico comportamento?
\end{ddtacheck}

\subsection{Essere verificabile non significa essere un requisito
separato}

Un FR deve descrivere un comportamento che sia possibile osservare e
verificare.

Questo non significa però che ogni parte testabile debba diventare un
Requirement distinto.

Due condizioni possono avere test diversi ma continuare ad appartenere
alla stessa obbligazione.

\textbf{Esempio semplice}

Supponiamo che un FR richieda di produrre una descrizione strutturata
composta da cinque campi.

Possiamo verificare separatamente la presenza e il contenuto dei cinque
campi, ma questo non significa che il progetto governi cinque
obbligazioni autonome.

Lo split avrebbe senso soltanto se uno o più di quei campi possedessero
davvero significato, ownership o ciclo di cambiamento indipendenti.
%</PAGE-CONTENT:33>

\clearpage
%<PAGE-DESC:34>FunctionalRequirement - non duplicazione e copertura della Decision
%<PAGE-CONTENT:34>
\PageHashFooter{\PageMDfive{34}}
\subsection{Evitare la duplicazione tra Decision e
FR}

Una Decision non dovrebbe contenere così tanto comportamento operativo
da lasciare al FunctionalRequirement soltanto una parafrasi.

Allo stesso modo, non dobbiamo rendere il FR deliberatamente vago solo
per farlo apparire diverso dal parent.

Il confine deve emergere dal significato:

\begin{itemize}
\tightlist
\item
  la Decision governa la scelta;
\item
  l\textquotesingle FR governa il comportamento necessario sotto quella
  scelta.
\end{itemize}

\textbf{Domanda di controllo}

\begin{ddtacheck}[title={Domanda di review}]
Se elimino questo candidato FR lasciando invariata la Decision, quale
comportamento richiesto dal progetto non è più vincolato?
\end{ddtacheck}

Se non riusciamo a identificare alcun comportamento, dobbiamo
riesaminare la necessità o la formulazione del candidato.

\subsection{Controllare che gli FR coprano la
Decision}

Dopo aver revisionato i singoli FR bisogna osservare anche la famiglia
nel suo insieme.

Immaginiamo che tutti gli FR figli della Decision siano soddisfatti.

\begin{ddtacheck}[title={Domanda di review}]
Esiste ancora un comportamento governato necessario alla Decision che
nessuno di questi FR copre?
\end{ddtacheck}

Se sì, possono esserci diverse spiegazioni:

\begin{itemize}
\tightlist
\item
  manca un FR;
\item
  uno degli FR è formulato in modo troppo astratto;
\item
  il comportamento è stato attribuito all\textquotesingle owner
  sbagliato;
\item
  la documentazione non contiene l\textquotesingle informazione
  necessaria.
\end{itemize}

Questo controllo serve a localizzare il problema. Non autorizza a
inventare il significato mancante.
%</PAGE-CONTENT:34>

\clearpage
%<PAGE-DESC:35>FunctionalRequirement - downstream utility regression e DermaTriage branch A
%<PAGE-CONTENT:35>
\PageHashFooter{\PageMDfive{35}}
\subsection{Controllare l'utilita' downstream dei branch sibling}

I controlli locali restano necessari ma non sufficienti. Una Decision può esprimere una scelta reale e il suo FR può essere formalmente corretto; tuttavia il branch può risultare inutile come \textbf{unità autonoma} se, per diventare implementabile, verificabile o analizzabile, deve essere ricongiunto integralmente a un branch sibling.

\begin{ddtacheck}[title={Branch-level Downstream Utility Test}]
Dopo aver assemblato gli FR, confrontare il branch completo con i sibling dello stesso MR e chiedere:
\begin{quote}
Se elimino questo branch come unità autonoma e trasferisco il suo significato nel branch che possiede i comportamenti che lo utilizzano, perdo una distinzione necessaria per implementazione, verifica, modifica indipendente o analisi?
\end{quote}
\textbf{YES}: il branch resta autonomo. \textbf{NO}: riaprire \texttt{MERGE}/\texttt{REWORK}. Il test non chiede se la Decision sia ``vera'', ma se la sua autonomia documentale produca valore downstream.
\end{ddtacheck}

\begin{ddtaexample}[title={DermaTriage - worked example completo, parte 1/2: branch semanticamente valido ma sospetto}]
\small
I due branch hanno lo stesso parent: \textbf{MR-01 - Valutazione di triage del caso dermatologico}.

\textbf{Branch A - DEC-MR01-02, Scala di priorita' P1--P4.}\par
\textbf{Decision}: ``DermaTriage usa la scala P1-P4 per rappresentare la priorita' operativa di triage.''\par
\textbf{Consequence}: la priorita' operativa e' uno dei livelli P1, P2, P3 o P4; le regole per assegnare il livello sono definite a parte.

\textbf{FR-MR01-02-01 - Dominio della priorita' P1--P4.}\par
\textbf{Functional obligation}: ogni priorita' operativa e' rappresentata mediante P1, P2, P3 o P4.\par
\textbf{Normative clause}: quando viene prodotta una priorita', il valore MUST appartenere a \{P1, P2, P3, P4\}.

La Decision non e' falsa: governa una scelta reale di dominio. Il problema emerge downstream: il suo FR ripete quel dominio ma non introduce un comportamento autonomo che possa essere implementato, verificato o analizzato senza il branch che determina concretamente la priorita'.
\end{ddtaexample}
%</PAGE-CONTENT:35>

\clearpage
%<PAGE-DESC:36>FunctionalRequirement - DermaTriage downstream utility e consolidazione pratica
%<PAGE-CONTENT:36>
\PageHashFooter{\PageMDfive{36}}
\begin{ddtaexample}[title={DermaTriage - worked example completo, parte 2/2: branch operativo e controfattuale}]
\small
\textbf{Branch B - DEC-MR01-04, separazione tra urgenza analitica e priorita' operativa.}\par
\textbf{Decision}: DermaTriage mantiene distinta la valutazione analitica dell'urgenza dalla priorita' operativa di triage; quest'ultima e' determinata successivamente e rappresentata mediante la scala P1--P4.\par
\textbf{Consequence}: la priorita' non e' l'esito analitico stesso: e' un risultato operativo distinto, appartenente al dominio P1--P4, derivato dall'urgenza e dalle ulteriori informazioni governate applicabili.

\textbf{FR-MR01-04-01 - Mapping urgenza/confidence verso P1--P4.}\par
\begin{center}
\begin{tabular}{ll}
\toprule
\textbf{Condizione} & \textbf{Priorita' richiesta} \\
\midrule
HIGH + confidence $>$ 0.85 & MUST produrre P1 \\
HIGH, negli altri casi applicabili & MUST produrre P2 \\
MEDIUM & MUST produrre P3 \\
LOW & MUST produrre P4 \\
\bottomrule
\end{tabular}
\end{center}

\textbf{Downstream Utility Test sul Branch A.} Rimuovendo DEC-MR01-02 e FR-MR01-02-01 come ramo autonomo e trasferendo la scelta P1--P4 nel contratto di DEC-MR01-04, non si perde alcuna distinzione necessaria per codice, test, modifica indipendente o analisi. Il dominio resta governato e il comportamento che lo rende effettivo resta nell'FR di mapping. Il Branch A era quindi \textbf{semanticamente valido ma strutturalmente superfluo}.
\end{ddtaexample}

\begin{ddtacore}[title={Consolidazione pratica risultante}]
La Decision mantenuta governa un solo contratto operativo coerente: trasformare un esito analitico in una \textbf{priorita' operativa distinta} appartenente al dominio P1--P4. Il FunctionalRequirement definisce la funzione concreta che associa urgenza/confidence ai quattro livelli. DEC-MR01-02 non viene eliminata perché ``sbagliata'', ma perché la sua autonomia non aggiunge utilita' downstream; FR-MR01-02-01 non giustifica il branch perché ripete soltanto il codominio già necessario al mapping.
\end{ddtacore}

\begin{ddtacheck}[title={Conservazione prima di accettare merge o rimozione}]
Prima della consolidazione verificare che restino rappresentati commitment, comportamenti, dominio, condizioni e binding materiali e che i gap restino visibili. Qui restano aperte, per esempio, l'applicabilita' del mapping al percorso symptom-only e la semantica degli eventuali tempi/SLA associati a P1--P4.
\end{ddtacheck}
%</PAGE-CONTENT:36>

\clearpage
%<PAGE-DESC:37>FunctionalRequirement - lacune documentali ed esempio DermaTriage
%<PAGE-CONTENT:37>
\PageHashFooter{\PageMDfive{37}}
\subsection{Una lacuna nella documentazione non è un
requisito}

Quando per comprendere il comportamento completo serve
un\textquotesingle informazione che le fonti non stabiliscono, quella
mancanza deve essere resa esplicita.

Non dobbiamo creare un FR per nascondere la lacuna o per rendere
l\textquotesingle albero apparentemente completo.

\textbf{Esempio DermaTriage}

Durante la review di MR-01 abbiamo trovato tre significati documentati:

\begin{itemize}
\tightlist
\item
  il triage può continuare anche senza immagine usando le informazioni
  sintomatologiche disponibili;
\item
  la priorità operativa appartiene al dominio P1-P4;
\item
  è documentata una regola che usa urgenza e confidence per selezionare
  un valore P1-P4 in un determinato percorso.
\end{itemize}

Le fonti analizzate non stabiliscono però che l\textquotesingle urgenza
prodotta dal percorso symptom-only debba attraversare necessariamente
quella stessa regola.

La risposta corretta non è inventare un ottavo FR.

La risposta corretta è registrare la domanda aperta:

\begin{ddtacheck}[title={Domanda di review}]
Come viene determinata la priorità P1-P4 quando
l\textquotesingle urgenza proviene dal percorso symptom-only?
\end{ddtacheck}

Finché le fonti non rispondono, il significato resta una \textbf{lacuna
documentale aperta}.
%</PAGE-CONTENT:37>

\clearpage
%<PAGE-DESC:38>FunctionalRequirement - obbligazione funzionale e dettagli della realizzazione
%<PAGE-CONTENT:38>
\PageHashFooter{\PageMDfive{38}}
\subsection{Obbligazione funzionale e dettagli della
realizzazione}

Questa distinzione è delicata perché non vogliamo né confondere il
requisito con l\textquotesingle implementazione né perdere informazioni
utili sul sistema reale.

Per ogni informazione chiediamoci quindi due cose diverse:

\begin{enumerate}
\def\labelenumi{\arabic{enumi}.}
\tightlist
\item
  \textbf{Qual è il comportamento che il progetto obbliga a garantire?}
\item
  \textbf{Come viene realizzato oggi quel comportamento?}
\end{enumerate}

Le due risposte possono cambiare in momenti diversi.

Un FR descrive anzitutto il comportamento governato. La documentazione
deve però conservare anche ciò che sappiamo della realizzazione
corrente: modelli, componenti, tecnologie, parametri, endpoint, formati,
configurazioni e binding.

\textbf{Separare non significa cancellare.}

Se, per esempio, un comportamento viene oggi realizzato da un
particolare modello ML, il nome del modello non deve scomparire dalla
documentazione solo perché non definisce necessariamente
l\textquotesingle identità del FunctionalRequirement.

Dobbiamo invece capire quale ruolo ha quel dettaglio.

\begin{itemize}
\tightlist
\item
  Se le fonti stabiliscono che il progetto \textbf{impone proprio quel
  dettaglio} come parte dell\textquotesingle obbligazione, esso può
  entrare nella semantica normativa dell\textquotesingle FR.
\item
  Se il dettaglio descrive \textbf{come l\textquotesingle obbligazione è
  realizzata oggi}, deve restare documentato e collegato
  all\textquotesingle FR, ma non diventa automaticamente parte
  dell\textquotesingle identità del requisito.
\item
  Se le fonti contengono il dettaglio ma non è ancora chiaro quale sia
  il suo ruolo normativo, lo manteniamo visibile e lo segnaliamo come
  punto da chiarire.
\end{itemize}

La domanda quindi non è:

\begin{ddtacheck}[title={Domanda di review}]
"Questo dettaglio è troppo tecnico e va eliminato?"
\end{ddtacheck}

ma:

\begin{ddtacheck}[title={Domanda di review}]
"Che cosa rappresenta questo dettaglio e dove deve essere governato
affinché non vada perso?"
\end{ddtacheck}

Resta inoltre una questione metodologica da affrontare in seguito:
alcuni vincoli o proprietà potrebbero richiedere un proprio tipo di
requisito specializzato, come già avviene per i
\passthrough{\lstinline!SecurityRequirement!} nel caso della sicurezza.

Questa sezione \textbf{non decide ancora} quali altri tipi debbano
esistere. Fissa soltanto un principio: l\textquotesingle informazione
non va eliminata perché non appartiene direttamente
all\textquotesingle identità dell\textquotesingle FR.
%</PAGE-CONTENT:38>

\clearpage
%<PAGE-DESC:39>FunctionalRequirement - information preservation e change impact
%<PAGE-CONTENT:39>
\PageHashFooter{\PageMDfive{39}}
\subsection{Nessuna informazione supportata dalle fonti deve andare
persa}

L\textquotesingle astrazione e la neutralizzazione servono a capire
\textbf{dove appartiene un\textquotesingle informazione}, non a
cancellarla.

Quando incontriamo un fatto supportato dalle fonti possiamo trovarci in
quattro situazioni.

\subsubsection{1. Il fatto appartiene direttamente
all\textquotesingle obbligazione
funzionale}

In questo caso può essere espresso nelle clausole normative
dell\textquotesingle FR.

\subsubsection{2. Il fatto appartiene a un altro owner o a
un\textquotesingle altra fase del ciclo di
vita}

L\textquotesingle informazione deve essere conservata presso
quell\textquotesingle owner e collegata al comportamento a cui si
riferisce.

Per esempio, un parametro di training può riguardare
l\textquotesingle evoluzione del modello e non il comportamento runtime
governato dallo specifico FR.

\subsubsection{3. Il fatto è documentato, ma non sappiamo ancora se sia
normativo}

Non lo promuoviamo automaticamente a requisito, ma non lo eliminiamo.
Rimane visibile come informazione supportata dalle fonti in attesa di
una classificazione migliore.

\subsubsection{4. L\textquotesingle informazione necessaria non è
documentata}

Non possiamo inventarla. Registriamo invece una lacuna esplicita.

In pratica, durante la review non dobbiamo chiederci soltanto:

\begin{ddtacheck}[title={Domanda di review}]
"Questa informazione appartiene all\textquotesingle FR?"
\end{ddtacheck}

Dobbiamo anche chiederci:

\begin{ddtacheck}[title={Domanda di review}]
"Se non appartiene all\textquotesingle FR, dove la documentiamo e come
manteniamo il collegamento con il comportamento a cui si riferisce?"
\end{ddtacheck}

\textbf{Quando qualcosa cambia}

Se cambia soltanto la realizzazione, la configurazione o un binding,
aggiorniamo la parte di documentazione che descrive quel dettaglio e la
relativa tracciabilità.

Se cambia invece il comportamento obbligatorio governato dal progetto,
riesaminiamo anche il FunctionalRequirement.

Se una tecnologia concreta è stata resa esplicitamente parte
dell\textquotesingle obbligazione dalle fonti, allora anche la sua
sostituzione può richiedere una modifica dell\textquotesingle FR.
%</PAGE-CONTENT:39>

\clearpage
%<PAGE-DESC:40>FunctionalRequirement - contratto L1 e presentazione L2
%<PAGE-CONTENT:40>
\PageHashFooter{\PageMDfive{40}}
\subsection{Contratto L1 minimo}

Il metamodel già governato resta la base. Questa revisione non introduce
nuovi campi nel FunctionalRequirement.

\begin{lstlisting}[basicstyle=\ttfamily\small]
Requirement [abstract]
    normativeClause : NormativeClause [1..*]

FunctionalRequirement extends Requirement
    parentDecision : Decision [1]
\end{lstlisting}

Questo significa che ogni FR:

\begin{itemize}
\tightlist
\item
  è un \passthrough{\lstinline!Requirement!};
\item
  possiede una o più \passthrough{\lstinline!normativeClause!};
\item
  ha esattamente una \passthrough{\lstinline!Decision!} parent.
\end{itemize}

Le informazioni aggiuntive utilizzate per rendere la documentazione
leggibile non devono diventare nuovi elementi del metamodel senza una
decisione metodologica esplicita.

\subsection{Come presentare un FR in modo
leggibile}

La rappresentazione documentale può mostrare in modo leggibile:

\begin{itemize}
\tightlist
\item
  ID;
\item
  stato di lifecycle;
\item
  authority;
\item
  tipo \passthrough{\lstinline!FunctionalRequirement!};
\item
  titolo;
\item
  Decision parent;
\item
  MacroRequirement derivato dal parent, quando aiuta la lettura;
\item
  una sintesi dell\textquotesingle obbligazione funzionale;
\item
  una o più \passthrough{\lstinline!normativeClause!};
\item
  riferimenti strutturati ed evidenze già governate;
\item
  eventuali \passthrough{\lstinline!SpecializedRequirement!} realmente
  stabiliti.
\end{itemize}

Questi elementi aiutano una persona a capire il requisito, ma non devono
creare una seconda semantica parallela rispetto al metamodel.

Informazioni come trigger, input, output, criteri di verifica,
realizzazione corrente e domande aperte possono essere molto utili
durante l\textquotesingle authoring e la review.

La loro utilità, però, non basta a trasformarle automaticamente in nuovi
campi L1.

Devono essere documentate usando le strutture e i meccanismi di
tracciabilità appropriati.
%</PAGE-CONTENT:40>

\clearpage
%<PAGE-DESC:41>FunctionalRequirement - scrittura delle clausole normative
%<PAGE-CONTENT:41>
\PageHashFooter{\PageMDfive{41}}
\subsection{Come scrivere le clausole
normative}

Le clausole normative devono rendere chiaro che cosa il progetto
richiede.

Per gli obblighi usiamo \passthrough{\lstinline!MUST!} e
\passthrough{\lstinline!MUST NOT!}.

\passthrough{\lstinline!SHOULD!} e \passthrough{\lstinline!MAY!} hanno
significati diversi:

\begin{itemize}
\tightlist
\item
  \passthrough{\lstinline!SHOULD!} esprime una raccomandazione
  governata;
\item
  \passthrough{\lstinline!MAY!} esprime un permesso governato.
\end{itemize}

Non vanno usati per attenuare un obbligo solo perché non siamo sicuri di
ciò che la fonte intende.

Se non sappiamo se qualcosa sia obbligatorio, il problema è
nell\textquotesingle authority o nella documentazione e deve rimanere
visibile.

La descrizione leggibile della \textbf{Functional obligation} aiuta il
lettore a comprendere l\textquotesingle FR, ma la fonte normativa
primaria del Requirement rimane l\textquotesingle insieme coerente delle
sue \passthrough{\lstinline!normativeClause!}.
%</PAGE-CONTENT:41>

\clearpage
%<PAGE-DESC:42>FunctionalRequirement - controllo finale del candidato
%<PAGE-CONTENT:42>
\PageHashFooter{\PageMDfive{42}}
\subsection{Controllo finale del candidato
FR}

Prima di accettare un candidato, prova a rispondere con evidenze
esplicite a queste domande:

\begin{enumerate}
\def\labelenumi{\arabic{enumi}.}
\tightlist
\item
  Quale parte delle fonti sostiene questo comportamento?
\item
  Quale unica Decision lo possiede?
\item
  Sto descrivendo un comportamento operativo oppure sto ripetendo la
  Decision?
\item
  Se elimino l\textquotesingle FR, quale comportamento richiesto smette
  di essere governato?
\item
  L\textquotesingle obbligazione possiede sufficiente autonomia da
  giustificare una propria identità?
\item
  L\textquotesingle eventuale split separa davvero due obbligazioni
  oppure frammenta artificialmente un unico comportamento?
\item
  Sto trasformando una tecnologia, un parametro, una rappresentazione o
  un binding in obbligo normativo senza sufficiente supporto?
\item
  Ogni valore concreto inserito nelle clausole normative è realmente
  sostenuto dalle fonti?
\item
  L\textquotesingle insieme degli FR copre il comportamento governato
  dalla Decision?
\item
  Le informazioni mancanti rimangono esplicite oppure le abbiamo
  completate per supposizione?
\end{enumerate}

Al termine della review il candidato può essere:

\begin{itemize}
\tightlist
\item
  mantenuto (\passthrough{\lstinline!KEEP!});
\item
  unito a un altro (\passthrough{\lstinline!MERGE!});
\item
  separato (\passthrough{\lstinline!SPLIT!});
\item
  riclassificato a un livello più basso
  (\passthrough{\lstinline!LOWER\_LEVEL!});
\item
  rielaborato (\passthrough{\lstinline!REWORK!});
\item
  lasciato in sospeso (\passthrough{\lstinline!HOLD!}).
\end{itemize}

Il codice sintetizza l\textquotesingle esito. La motivazione e
l\textquotesingle evidenza spiegano perché quell\textquotesingle esito è
corretto.

Nessuna decisione deve dipendere dal raggiungimento di un numero atteso
di FR.

\begin{ddtacheck}[title={Controllo obbligatorio - percorso dati e contratti informativi}]
Dopo aver identificato i FunctionalRequirement di una Decision, enumerare ogni informazione prodotta, consumata, resa disponibile o persistita. Per ogni coppia producer/consumer verificare se esiste un transfer semanticamente rilevante. Ogni transfer candidato deve poter rendere identificabili source, destination e content; quando il segmento deve essere analizzato, qualificato, realizzato o protetto autonomamente, deve poter ricevere identita'. Per ogni contenuto trasferito deve essere verificata la sufficienza del relativo information/data contract. Informazioni necessarie ma non stabilite dalla fonte devono restare esplicitamente \texttt{NOT SPECIFIED} / diagnostic gap e non essere ricostruite per plausibilita'.

Questo controllo non introduce nuovi campi L1 e non impone un FunctionalRequirement per ogni scambio: transfer, contratto informativo, binding e realization devono ancora essere classificati secondo source support, semantic owner e livello corretto.
\end{ddtacheck}
%</PAGE-CONTENT:42>

\clearpage
%<PAGE-DESC:43>FunctionalRequirement - controllo famiglia, downstream utility e analisi indipendente
%<PAGE-CONTENT:43>
\PageHashFooter{\PageMDfive{43}}
\subsection{Controllare la famiglia completa degli FR}

Anche quando ogni FR sembra corretto individualmente, la famiglia può avere problemi. Dopo la review dei singoli requisiti verifica quindi che:

\begin{itemize}
\tightlist
\item ogni FR abbia una sola Decision parent;
\item due FR non governino inutilmente lo stesso comportamento;
\item un unico comportamento non sia stato frammentato senza motivo;
\item nessun significato operativo governato sia rimasto senza owner;
\item i dettagli di realizzazione, configurazione e binding siano ancora presenti e tracciabili;
\item le lacune ancora aperte siano visibili;
\item eventuali modifiche alle Decision non abbiano reso obsoleti i relativi FR;
\item ogni branch Decision ancora autonomo abbia superato il \textbf{Downstream Utility Test}.
\end{itemize}

\begin{ddtacore}[title={Feedback obbligatorio verso la Decision family}]
Se una Decision non possiede FR propri, oppure i suoi FR non introducono alcun comportamento downstream distinto, non inventare requisiti e non chiudere automaticamente il ramo. Riportare il branch alla Decision-family review e verificare se la sua autonomia serve davvero a implementazione, verifica, modifica indipendente o analisi. In caso contrario valutare \texttt{MERGE} o \texttt{REWORK} preservando il significato nel branch operativo corretto.
\end{ddtacore}

\textbf{Il numero degli FR è un risultato, non un obiettivo.} Una famiglia può contenerne più o meno rispetto a una versione precedente se le fonti o la comprensione del significato lo giustificano.

\subsection{Confrontare la derivazione con un\textquotesingle analisi indipendente}
Una guida è più utile se persone diverse, partendo dalle stesse fonti, ricostruiscono significati compatibili. Possiamo quindi chiedere a un altro ricercatore di derivare gli FR indipendentemente e confrontare i risultati. Lo scopo non è stabilire chi ``ha vinto'', ma individuare:

\begin{itemize}
\tightlist
\item punti ambigui nelle fonti;
\item passaggi poco chiari nella guida;
\item criteri di split/merge interpretati diversamente;
\item informazioni tecniche collocate in livelli differenti;
\item eventuali esempi della guida che hanno orientato troppo la risposta.
\end{itemize}

L\textquotesingle analisi esterna e l\textquotesingle output di un LLM restano \textbf{evidenze di interpretazione}, non authority di progetto. La maggioranza non decide il significato: per riconciliare le divergenze si torna alle fonti, all'owner del significato e ai controlli previsti dalla guida.
%</PAGE-CONTENT:43>

\clearpage
%<PAGE-DESC:44>FunctionalRequirement - classificazione delle differenze tra analisi
%<PAGE-CONTENT:44>
\PageHashFooter{\PageMDfive{44}}
\subsection{Come classificare le differenze tra due
analisi}

Quando due derivazioni non coincidono, conviene descrivere \textbf{su
che cosa} stanno divergendo prima di tentare di risolvere il problema.

La tassonomia usata nel lavoro sperimentale comprende:

\begin{itemize}
\tightlist
\item
  \textbf{split o merge (\passthrough{\lstinline!SPLIT\_MERGE!})}: gli
  analisti non concordano su quante obbligazioni autonome esistano;
\item
  \textbf{authority della fonte
  (\passthrough{\lstinline!SOURCE\_AUTHORITY!})}: non concordano su
  quale fonte autorizzi un significato;
\item
  \textbf{Decision o FR (\passthrough{\lstinline!FR\_VS\_DECISION!})}:
  lo stesso significato viene collocato a livelli diversi;
\item
  \textbf{FR o realizzazione
  (\passthrough{\lstinline!FR\_VS\_REALIZATION!})}: un dettaglio viene
  interpretato da uno come requisito e dall\textquotesingle altro come
  realizzazione;
\item
  \textbf{FR o binding (\passthrough{\lstinline!FR\_VS\_BINDING!})}:
  diverge la classificazione di un collegamento concreto;
\item
  \textbf{FR o parametro (\passthrough{\lstinline!FR\_VS\_PARAMETER!})}:
  diverge il ruolo di un valore o parametro;
\item
  \textbf{FR o specializzazione
  (\passthrough{\lstinline!FR\_VS\_SPECIALIZATION!})}: non è chiaro se
  il significato appartenga alla funzione ordinaria o a una proprietà
  specializzata;
\item
  \textbf{ownership (\passthrough{\lstinline!OWNERSHIP!})}: non
  c\textquotesingle è accordo su chi debba governare il significato;
\item
  \textbf{significato mancante
  (\passthrough{\lstinline!MISSING\_SEMANTICS!})}: una delle analisi
  evidenzia informazioni necessarie ma non documentate;
\item
  \textbf{ambiguità della guida
  (\passthrough{\lstinline!GUIDE\_AMBIGUITY!})}: la stessa regola
  metodologica permette letture differenti;
\item
  \textbf{influenza degli esempi
  (\passthrough{\lstinline!GUIDE\_EXAMPLE\_ANCHORING!})}: un esempio
  della guida può avere suggerito la soluzione;
\item
  \textbf{differenza solo rappresentativa
  (\passthrough{\lstinline!REPRESENTATIONAL\_ONLY!})}: cambia la forma
  ma non il significato;
\item
  \textbf{altro (\passthrough{\lstinline!OTHER!})}: differenza non
  coperta dalle categorie precedenti.
\end{itemize}

Queste etichette non risolvono la divergenza. Servono a renderla
osservabile e confrontabile.
%</PAGE-CONTENT:44>

\clearpage
%<PAGE-DESC:45>FunctionalRequirement - replica indipendente e supporto LLM
%<PAGE-CONTENT:45>
\PageHashFooter{\PageMDfive{45}}
\subsection{Preparare una replica realmente
indipendente}

Se vogliamo capire quanto la guida sia ripetibile, il ricercatore
indipendente non deve conoscere in anticipo la soluzione che vogliamo
verificare.

Nel pacchetto di lavoro vanno quindi evitati, per quanto possibile:

\begin{itemize}
\tightlist
\item
  vecchi ID degli FR del progetto target;
\item
  numero atteso di requisiti;
\item
  esempi già risolti sullo stesso progetto;
\item
  commenti che rivelano quale split o merge sia stato accettato;
\item
  gap formulati in modo tale da suggerire già la classificazione finale.
\end{itemize}

Il ricercatore deve poter arrivare alla propria derivazione usando le
fonti e le regole metodologiche disponibili.

Una sequenza utile è:

\begin{lstlisting}[basicstyle=\ttfamily\small]
derivazione indipendente
        |
confronto dei significati
        |
classificazione delle differenze
        |
riconciliazione umana tornando alle fonti
        |
minima correzione necessaria della guida
        |
nuovo controllo sullo stesso ramo
        |
nuova replica indipendente senza suggerimenti sulla risposta
\end{lstlisting}

In questo modo le differenze tra le analisi diventano uno strumento per
capire \textbf{dove la metodologia o la documentazione sono ancora
ambigue}.

\subsection{Uso di un LLM come supporto
all\textquotesingle authoring}

Un LLM può aiutare a cercare candidati, organizzare evidenze e applicare
in modo ripetibile le domande della guida.

Non può però creare significato di progetto.

Il modello deve quindi seguire gli stessi principi
dell\textquotesingle autore umano:

\begin{itemize}
\tightlist
\item
  usare soltanto le fonti autorizzate per stabilire il significato;
\item
  non trasformare automaticamente i dettagli tecnici in requisiti;
\item
  non eliminare dettagli documentati solo perché sono classificati come
  realizzazione o configurazione;
\item
  non usare vecchi FR come prova che una nuova derivazione sia corretta;
\item
  non colmare informazioni mancanti con conoscenza di dominio o con una
  soluzione plausibile;
\item
  rendere visibili dubbi e lacune invece di nasconderli.
\end{itemize}

Il prompt canonico per i FunctionalRequirement rimane uno strumento
operativo separato dalla spiegazione metodologica.

La guida deve spiegare \textbf{come ragionare}; il prompt deve aiutare
il modello ad applicare quel ragionamento in modo controllato.
%</PAGE-CONTENT:45>

\clearpage
%<PAGE-DESC:46>Indice di integrita' delle pagine
%<PAGE-CONTENT:46>
\PageHashFooter{\PageMDfive{46}}
\section*{Indice di integrita' delle pagine}
\addcontentsline{toc}{section}{Indice di integrita' delle pagine}
\begin{ddtacore}[title={Regola di verifica}]
Ogni pagina possiede un MD5 del proprio contenuto canonico. Il calcolo usa il contenuto LaTeX compreso tra i marker \texttt{PAGE-CONTENT}; il valore MD5 stampato sulla pagina e' escluso dal payload della pagina stessa. Nell'indice, gli hash delle altre pagine sono risolti prima del calcolo, mentre l'hash della pagina indice viene sostituito dal token canonico \texttt{<SELF>} per evitare una dipendenza circolare.
\end{ddtacore}
\vspace{2mm}
\begingroup
\tiny
\renewcommand{\arraystretch}{0.64}
\setlength{\tabcolsep}{2.2pt}
\renewcommand{\IntegrityRow}[3]{#1 & #2 & {\tiny\ttfamily #3} \\}
\begin{longtable}{L{8mm}L{89mm}L{47mm}}
\toprule
\textbf{Pag.} & \textbf{Contenuto} & \textbf{MD5 contenuto} \\
\midrule
%<PAGE-INTEGRITY-ROWS>
\IntegrityRow{1}{Frontespizio della nuova guida di authoring}{\PageMDfive{1}}
\IntegrityRow{2}{Project Problem Framing - definizione e regole di authoring}{\PageMDfive{2}}
\IntegrityRow{3}{Project Problem Framing - procedura umana, gate ed esempio DermaTriage}{\PageMDfive{3}}
\IntegrityRow{4}{Project Problem Framing - LLM Execution Profile e prompt canonico}{\PageMDfive{4}}
\IntegrityRow{5}{Project Problem Framing - regola LLM e classificazione dei campi}{\PageMDfive{5}}
\IntegrityRow{6}{Project Problem Framing - formato esatto dell'output LLM}{\PageMDfive{6}}
\IntegrityRow{7}{Project Problem Framing - gate LLM e nota sui controlli quantitativi}{\PageMDfive{7}}
\IntegrityRow{8}{Uso della metodologia - authoring nativo e reconstruction}{\PageMDfive{8}}
\IntegrityRow{9}{MacroRequirement - definizione e rapporto con framing e Decision}{\PageMDfive{9}}
\IntegrityRow{10}{MacroRequirement - procedura umana per individuare i candidati}{\PageMDfive{10}}
\IntegrityRow{11}{MacroRequirement - test di separazione, split e merge}{\PageMDfive{11}}
\IntegrityRow{12}{MacroRequirement - campi documentali e regole di compilazione}{\PageMDfive{12}}
\IntegrityRow{13}{MacroRequirement - template, split merge e STOP AT MR}{\PageMDfive{13}}
\IntegrityRow{14}{MacroRequirement - sufficienza semantica e diagnosi dei problemi}{\PageMDfive{14}}
\IntegrityRow{15}{MacroRequirement - profilo LLM e prompt canonico parte 1}{\PageMDfive{15}}
\IntegrityRow{16}{MacroRequirement - prompt canonico LLM parte 2}{\PageMDfive{16}}
\IntegrityRow{17}{MacroRequirement - prompt LLM campi e schema di analisi}{\PageMDfive{17}}
\IntegrityRow{18}{MacroRequirement - formato esatto output LLM e gate finale}{\PageMDfive{18}}
\IntegrityRow{19}{Decision - definizione confine e discovery}{\PageMDfive{19}}
\IntegrityRow{20}{Decision - worksheet candidate e source support}{\PageMDfive{20}}
\IntegrityRow{21}{Decision - MR stability neutralization semantic level}{\PageMDfive{21}}
\IntegrityRow{22}{Decision - independent governance e removal replacement}{\PageMDfive{22}}
\IntegrityRow{23}{Decision - decomposition, semantic owner e utilita' downstream}{\PageMDfive{23}}
\IntegrityRow{24}{Decision - esempio completo e MR prose leakage}{\PageMDfive{24}}
\IntegrityRow{25}{Decision - family regression, downstream utility e closure}{\PageMDfive{25}}
\IntegrityRow{26}{Decision - LLM execution profile prompt canonico parte 1}{\PageMDfive{26}}
\IntegrityRow{27}{Decision - prompt LLM candidate gates}{\PageMDfive{27}}
\IntegrityRow{28}{Decision - prompt LLM decomposition, downstream regression e output}{\PageMDfive{28}}
\IntegrityRow{29}{FunctionalRequirement - definizione, termini e domanda fondamentale}{\PageMDfive{29}}
\IntegrityRow{30}{FunctionalRequirement - confine Decision-FR, parent unico e classificazioni di review}{\PageMDfive{30}}
\IntegrityRow{31}{FunctionalRequirement - individuazione source-first dei comportamenti}{\PageMDfive{31}}
\IntegrityRow{32}{FunctionalRequirement - scheda candidate, branch utility e obbligazione coerente}{\PageMDfive{32}}
\IntegrityRow{33}{FunctionalRequirement - split e verificabilita' senza frammentazione}{\PageMDfive{33}}
\IntegrityRow{34}{FunctionalRequirement - non duplicazione e copertura della Decision}{\PageMDfive{34}}
\IntegrityRow{35}{FunctionalRequirement - downstream utility regression e DermaTriage branch A}{\PageMDfive{35}}
\IntegrityRow{36}{FunctionalRequirement - DermaTriage downstream utility e consolidazione pratica}{\PageMDfive{36}}
\IntegrityRow{37}{FunctionalRequirement - lacune documentali ed esempio DermaTriage}{\PageMDfive{37}}
\IntegrityRow{38}{FunctionalRequirement - obbligazione funzionale e dettagli della realizzazione}{\PageMDfive{38}}
\IntegrityRow{39}{FunctionalRequirement - information preservation e change impact}{\PageMDfive{39}}
\IntegrityRow{40}{FunctionalRequirement - contratto L1 e presentazione L2}{\PageMDfive{40}}
\IntegrityRow{41}{FunctionalRequirement - scrittura delle clausole normative}{\PageMDfive{41}}
\IntegrityRow{42}{FunctionalRequirement - controllo finale del candidato}{\PageMDfive{42}}
\IntegrityRow{43}{FunctionalRequirement - controllo famiglia, downstream utility e analisi indipendente}{\PageMDfive{43}}
\IntegrityRow{44}{FunctionalRequirement - classificazione delle differenze tra analisi}{\PageMDfive{44}}
\IntegrityRow{45}{FunctionalRequirement - replica indipendente e supporto LLM}{\PageMDfive{45}}
\IntegrityRow{46}{Indice di integrita' delle pagine}{\PageMDfive{46}}
%</PAGE-INTEGRITY-ROWS>
\bottomrule
\end{longtable}
\endgroup
\vspace{2mm}
{\small\color{DDTAGray}
A ogni passo, una modifica intenzionale cambia l'MD5 della pagina interessata. Una variazione inattesa dell'MD5 di una pagina gia' approvata impone la revisione della modifica prima di proseguire.}
%</PAGE-CONTENT:46>
\end{document}
